\documentclass[11pt]{article}

\usepackage[margin=1in]{geometry}
\usepackage{amsmath,amssymb,amsfonts,mathtools,amsthm,mathrsfs}
\usepackage{bm}
\usepackage{booktabs,array,longtable,tabularx}
\usepackage{enumitem}
\usepackage{xcolor}
\usepackage{xurl}
\usepackage[hidelinks]{hyperref}

\numberwithin{equation}{section}
\setlist{nosep,leftmargin=*}
\setlength{\emergencystretch}{2em}

\newcommand{\CP}{\mathbb{CP}}
\newcommand{\CC}{\mathbb C}
\newcommand{\RR}{\mathbb R}
\newcommand{\ZZ}{\mathbb Z}
\newcommand{\cO}{\mathcal O}
\newcommand{\cH}{\mathcal H}
\newcommand{\sltwo}{\mathfrak{sl}_2(\mathbb C)}
\newcommand{\Ktr}{K_{\mathrm{tr}}}
\newcommand{\GeV}{\,\mathrm{GeV}}
\newcommand{\dd}{\mathrm d}
\newcommand{\ii}{\mathrm i}
\newcommand{\Tr}{\operatorname{Tr}}
\newcommand{\Ind}{\operatorname{Ind}}
\newcommand{\ind}{\operatorname{ind}}
\newcommand{\Arf}{\operatorname{Arf}}
\newcommand{\rank}{\operatorname{rank}}
\newcommand{\divisor}{\operatorname{div}}
\newcommand{\Span}{\operatorname{span}}
\newcommand{\status}[1]{\textsf{#1}}

\definecolor{darkgreen}{RGB}{0,105,55}
\definecolor{darkred}{RGB}{150,25,25}
\definecolor{darkblue}{RGB}{20,65,140}

\newtheorem{theorem}{Theorem}[section]
\newtheorem{proposition}[theorem]{Proposition}
\newtheorem{lemma}[theorem]{Lemma}
\newtheorem{corollary}[theorem]{Corollary}
\newtheorem{definition}[theorem]{Definition}
\newtheorem{assumption}[theorem]{Bridge axiom}
\newtheorem{remark}[theorem]{Remark}
\newtheorem{gate}[theorem]{Promotion gate}

\title{Another Physics and Route E:\
First-Principles Separation Theorems, an Exact
Discriminant--Killing Bridge, and a Reversible Research Ledger}
\author{Route F research note}
\date{17 July 2026}

\begin{document}
\maketitle

\begin{abstract}
This ledger asks which parts of the manuscript \emph{Another physics} can be
connected to Route E without confusing an algebraic resemblance with a
physical derivation.  The literal identification of local energy density with
time, phase-dependent disappearance of invariant mass, fundamental negative
energy, and indestructible entanglement is ruled out under ordinary local,
positive-energy relativistic quantum theory.  A conditional complex-scalar
completion does support finite-charge Q-balls, but its internal phase is not
itself an observable clock.

The strongest exact connection is new but elementary.  On Route E's
H3$^+$-conditional branch,
\(H^0(\CP^1,T_{\CP^1})\simeq H^0(\CP^1,\cO(2))\simeq\sltwo\).
For a binary quadratic, its discriminant simultaneously detects two distinct
zeros, a regular-semisimple M\"obius flow, a nonzero Killing norm, and a
nonzero Route-E second-transvectant contact.  Thus the two-center divisor and
\(\Ktr\) are not merely parallel observations: they are two readings of the
same invariant.  A Fubini--Study moment map can further label the two fixed
points by opposite bounded charges without introducing negative Hamiltonian
energy.

A candidate action-level bridge makes the Route-E messenger mixing depend on
a Q-ball field, \(\lambda(\Phi)=g\Phi/\Lambda\), giving
\(\zeta_{\rm eff}=g^2\Phi^2/\Lambda^2\) and hence the same weight-two phase
law as Route E.  This is explicitly a bridge axiom, not a promoted result: an
exact Q-ball \(U(1)\) makes every invariant observable stationary, while an
observable phase modulation requires relative-phase breaking and therefore
reopens charge leakage and soliton stability.  All claims are classified,
derived, and tied to a reproducible numerical ledger; no physics promotion is
authorized by the present note.
\end{abstract}

\tableofcontents

\section{Provenance, scope, and logical contract}

The source under review is the 15-page file
\texttt{Another physics.pdf}, supplied from the local Downloads directory,
with PDF metadata title
\texttt{Another physics\_20230910.docx} and SHA-256
\begin{center}
\small\texttt{7355cf2237513a6e943555ac2720871a5e81915fd770541baea9d87eccf2654e}.
\end{center}
The Route-E reference state for this ledger is repository commit
\begin{center}
\small\texttt{e7ba020227853a3b731235b5f3a42083559e7ef2}.
\end{center}
The governing Route-E sources are \path{route_E/tex/route_e_first_principles.tex},
\path{route_E/tex/three_families_killing_contact.tex}, and the canonical
dynamics registry \path{route_E/code_dyn/dyn_claim_registry.json}.

Four logical labels are used throughout.

\begin{longtable}{>{\raggedright\arraybackslash}p{0.19\textwidth}
                  >{\raggedright\arraybackslash}p{0.72\textwidth}}
\toprule
Label & Meaning \\
\midrule
\endfirsthead
\toprule
Label & Meaning \\
\midrule
\endhead
\status{theorem} & Consequence of explicitly stated mathematical or physical
assumptions, with a proof in this ledger or an identified Route-E proof.\\
\status{conditional completion} & A consistent model after a new action,
field content, or boundary condition has been chosen.  It does not derive
that choice from the source manuscript.\\
\status{bridge axiom} & An identification between objects belonging to two
different models.  Similar degree, phase, or dimension is not a proof of the
identification.\\
\status{numerical check} & Reproducible arithmetic within a declared
approximation.  It corroborates formulas but cannot promote a bridge axiom.\\
\bottomrule
\end{longtable}

\begin{gate}[Non-promotion]
The machine card may be green only with the non-promoting status
\path{mechanical_separation_checks_pass_no_physics_promotion}.  Its Boolean
field \path{physics_promotion_allowed} must remain \texttt{false}
until an action-level embedding, stability analysis, and all affected Route-E
phenomenology gates pass.
\end{gate}

\section{Covariant dictionary and literal-theory boundary}

The source concepts are useful only after observer, sign, and conservation
laws are made explicit.  We use metric signature \((-+++ )\).  For an
observer with unit timelike four-velocity \(u^\mu\), the local measured energy
density is
\begin{equation}
 E_d(x;u):=T_{\mu\nu}(x)u^\mu u^\nu.
 \label{eq:Ed}
\end{equation}
This is a scalar after \(u\) is specified, but it is not an observer-independent
replacement for spacetime or proper time.  A viable information variable is
instead a Noether charge, an entropy/correlation functional, or an order
parameter.  In the Q-ball completion below it is
\begin{equation}
 E_i\quad\rightsquigarrow\quad Q_I=\int_\Sigma \dd\Sigma_\mu\,j_I^\mu,
 \qquad \nabla_\mu j_I^\mu=0.
 \label{eq:EiQ}
\end{equation}
The symbols \(E_p,E_n\) will later be reinterpreted as bounded moment-map
polarities \(\mu=+1,-1\), not as positive and negative Hamiltonian energies.

\subsection{Local density cannot determine clock rate}

\begin{theorem}[Local-density time no-go]
No function of the pointwise scalar \(E_d(x;u)\) alone can reproduce the
gravitational clock rate in all solutions of general relativity.
\end{theorem}

\begin{proof}
Consider a static thin spherical shell of mass \(M\) and radius \(R>2GM/c^2\).
The vacuum interior is flat but its lapse is fixed by matching to the exterior:
\begin{equation}
 \dd s^2_{r<R}
 =-\left(1-\frac{2GM}{Rc^2}\right)c^2\dd t^2
 +\dd r^2+r^2\dd\Omega^2.
 \label{eq:shellmetric}
\end{equation}
Both the shell interior and spatial infinity have \(T_{\mu\nu}=0\), hence
\(E_d=0\), but stationary clocks obey
\begin{equation}
 \frac{\dd\tau_{\rm in}}{\dd t_\infty}
 =\sqrt{1-\frac{2GM}{Rc^2}}\ne1.
 \label{eq:shellclock}
\end{equation}
Two points with the same local density therefore have different clock rates.
The missing datum is the nonlocal metric solution/boundary condition, not a
different definition of local energy density.
\end{proof}

For Earth-scale numbers, \(GM/(Rc^2)=6.9613\times10^{-10}\), corresponding at
leading order to about \(60.15\,\mu\mathrm{s}\) per day.  In the laboratory
weak-field limit,
\begin{equation}
 \frac{\Delta\nu}{\nu}=\frac{g\Delta h}{c^2}
 =1.091\times10^{-18}\left(\frac{\Delta h}{1\,\mathrm{cm}}\right).
 \label{eq:clockweak}
\end{equation}
This scale has been directly tested by a centimeter atomic-clock network
\cite{Zheng2023} and is an immediate constraint on any replacement for metric
proper time.

The same local \(E_d\) can also curve spacetime differently because pressure
matters.  Tracing
\begin{equation}
 G_{\mu\nu}+\Lambda g_{\mu\nu}=\frac{8\pi G}{c^4}T_{\mu\nu}
\end{equation}
and using \(T=-\varepsilon+3p\) gives
\begin{equation}
 R=4\Lambda+\frac{8\pi G}{c^4}(\varepsilon-3p).
 \label{eq:RicciTrace}
\end{equation}
Dust and radiation may have the same \(\varepsilon\), but \(p=0\) and
\(p=\varepsilon/3\) give different Ricci scalars.  A density-only clock or
gravity law therefore discards essential stress data.

\subsection{Mass, energy sign, and entanglement}

\begin{theorem}[Phase-toggle no-go in a fixed particle representation]
For a stable one-particle irreducible representation of the Poincar\'e group,
the rest mass cannot alternate between \(m\ne0\) and \(m=0\) as an unobserved
internal phase changes while the state remains in that representation.
\end{theorem}

\begin{proof}
The mass is fixed by the Casimir
\begin{equation}
 P_\mu P^\mu=-m^2c^2.
 \label{eq:casimir}
\end{equation}
It commutes with the representation.  If a phase trajectory changed the
eigenvalue from \(-m^2c^2\) to zero, it would leave the massive irreducible
representation and enter a massless one.  A periodic \emph{coupling}, spectral
weight, quasiparticle gap, or detector response may vanish; invariant rest
mass in a fixed representation may not.
\end{proof}

Likewise, a literal independent negative-energy oscillator,
\begin{equation}
 H=\omega a^\dagger a-\Omega b^\dagger b,
\end{equation}
has \(\inf\operatorname{spec}H=-\infty\), since
\(\langle n_b|H|n_b\rangle=-n_b\Omega\).  A stable completion must reinterpret
the sign as charge, orientation, a hole relative to a filled state, or an
effective-medium variable; it cannot be a freely excitable fundamental
negative-energy sector.

Entanglement is not indestructible either.  For
\begin{equation}
 \rho_p=(1-p)|\Phi^+\rangle\langle\Phi^+|+\frac p4 I_4,
 \qquad |\Phi^+\rangle=\frac{|00\rangle+|11\rangle}{\sqrt2},
\end{equation}
the partial transpose eigenvalues are
\begin{equation}
 \operatorname{spec}(\rho_p^{T_B})
 =\left\{\frac{2-p}{4},\frac{2-p}{4},\frac{2-p}{4},
 \frac{-2+3p}{4}\right\}.
 \label{eq:bellpt}
\end{equation}
For \(p\ge2/3\) the state is PPT and, in \(2\times2\), separable.  What is
forbidden is superluminal signalling, not dynamical loss of entanglement.

\subsection{Other immediate separation results}

For a massive electron of finite Lorentz factor \(\gamma\),
\begin{equation}
 \beta:=\frac vc=\sqrt{1-\gamma^{-2}}<1.
\end{equation}
Magnetism is encoded in the gauge curvature
\(F_{\mu\nu}=\partial_\mu A_\nu-\partial_\nu A_\mu\), not in electrons moving
exactly at \(c\).  Event-horizon escape is similarly excluded by causal
structure; jets must be launched outside the horizon.

If a proposed galactic potential is parabolic, \(\Phi_N(r)=a r^2\), then
\begin{equation}
 \frac{v_c^2}{r}=\frac{\dd\Phi_N}{\dd r}=2ar
 \quad\Longrightarrow\quad v_c=\sqrt{2a}\,r,
\end{equation}
which is solid-body rotation, not an asymptotically flat curve.  A concrete
modified-gravity alternative must fit both dynamics and lensing; recent weak
lensing reaches far beyond classical rotation-curve radii \cite{Mistele2024}.
Finally a decay half-life alone supplies no force.  Only momentum flux does;
for perfectly absorbed radiation
\begin{equation}
 F=\frac Pc=3.3356\,\mathrm{nN}\left(\frac P{1\,\mathrm W}\right).
\end{equation}

\section{The Route-E theorem boundary imported into Route F}

Route E assumes a self-carried family space
\begin{equation}
 V=H^0(\Sigma,L),\qquad L\simeq T_\Sigma,
 \qquad V\simeq H^0(\Sigma,T_\Sigma)=\mathfrak{aut}(\Sigma).
 \label{eq:selfcarried}
\end{equation}
The original invariant-contact hypothesis H3 requires a nondegenerate
symmetric invariant bilinear form.  It gives only \(N_{\rm fam}\le3\), since
the genus-one one-dimensional Abelian algebra admits such a form.  The exact
selection \(N_{\rm fam}=3\) requires the explicit stronger hypothesis H3$^+$:
the contact is the nondegenerate adjoint-trace/Killing form.  On that selected
branch,
\begin{equation}
 \Sigma\simeq\CP^1,\qquad T_{\CP^1}\simeq\cO(2),\qquad
 V\simeq\sltwo,\qquad \dim_{\mathbb C}V=3.
 \label{eq:routeEbranch}
\end{equation}
Thus every result below that invokes this branch is conditional on H3$^+$;
the subsequent algebra needs no additional physical axiom once the branch is
chosen.

Riemann--Roch and Serre duality give the dimension transparently:
\begin{align}
 h^0(\cO(2))-h^1(\cO(2))&=\deg\cO(2)+1=3,\\
 h^1(\cO(2))&=h^0(K\otimes\cO(-2))=h^0(\cO(-4))=0,
\end{align}
so \(h^0=3\).  Here \(\cO(2)\) is a degree-two holomorphic line bundle, not
the orthogonal group \(O(2)\).

In Route E's spherical basis \((T_{+1},T_0,T_{-1})\),
\begin{equation}
 \Ktr=\frac1{\sqrt3}
 \begin{pmatrix}
 0&0&-1\\[2pt]
 0&1&0\\[2pt]
 -1&0&0
 \end{pmatrix},
 \qquad B=2\sqrt3\,\Ktr,
 \qquad \Ktr^2=\frac13 I.
 \label{eq:Ktr}
\end{equation}
This complex symmetric invariant pairing is not a Hermitian energy operator.
Its indefinite real-slice signs cannot be read as positive and negative
physical energies.

\section{Exact discriminant--two-center--Killing theorem}

\subsection{Global vector fields and the adjoint representation}

Let \(\xi=v/u\) on an affine chart of \(\CP^1\).  A holomorphic vector field
has the local form \(X=q(\xi)\partial_\xi\).  On the other chart
\(\eta=1/\xi\),
\begin{equation}
 q(\xi)\partial_\xi=-\eta^2q(1/\eta)\partial_\eta.
\end{equation}
Regularity at \(\eta=0\) therefore forces \(\deg q\le2\):
\begin{equation}
 X_q=(a+b\xi+c\xi^2)\partial_\xi.
 \label{eq:vectorfield}
\end{equation}
The basis
\begin{equation}
 E=\partial_\xi,\qquad H=-2\xi\partial_\xi,
 \qquad F=-\xi^2\partial_\xi
\end{equation}
obeys
\begin{equation}
 [H,E]=2E,\qquad[H,F]=-2F,\qquad[E,F]=H,
\end{equation}
which proves \(H^0(\CP^1,T)\simeq\sltwo\).

Homogenizing \eqref{eq:vectorfield} gives the binary quadratic
\begin{equation}
 p(u,v)=a u^2+buv+cv^2,\qquad \Delta(p)=b^2-4ac.
 \label{eq:binaryquad}
\end{equation}
The representation statement is precisely
\begin{equation}
 H^0(\CP^1,\cO(2))\simeq\operatorname{Sym}^2(\mathbb C^2)^*,
\end{equation}
and the invariant epsilon tensor identifies the dual with
\(\operatorname{Sym}^2\mathbb C^2\) for \(SL(2)\).

\begin{theorem}[Discriminant--Killing equivalence]
\label{thm:disc-killing}
For a nonzero \(X_q\in H^0(\CP^1,T)\), the following are equivalent:
\begin{enumerate}
\item its zero divisor is a sum of two distinct points;
\item \(\Delta(p)\ne0\);
\item the corresponding \(\sltwo\) element is regular semisimple;
\item its Killing norm is nonzero;
\item its Route-E contact norm \(\Ktr(X_q,X_q)\) is nonzero.
\end{enumerate}
The complement \(\Delta=0\) is the nonzero nilpotent/contact-null cone and
has one double zero.
\end{theorem}

\begin{proof}
Take the standard projective infinitesimal matrix
\begin{equation}
 A_q=
 \begin{pmatrix}-b/2&-c\\ a&b/2\end{pmatrix}.
\end{equation}
Direct multiplication gives
\begin{equation}
 A_q^2=\frac{\Delta(p)}4 I.
 \label{eq:AqSquare}
\end{equation}
Hence \(\Delta\ne0\) gives two nonzero opposite eigenvalues and a regular
semisimple element.  If \(\Delta=0\) and \(A_q\ne0\), then \(A_q^2=0\), so
the element is nilpotent.  The roots of \(p\) are distinct exactly when its
discriminant is nonzero; when it vanishes, the homogeneous quadratic is a
square and has a double projective root.  Counting multiplicity, the total
degree is always
\begin{equation}
 \deg\divisor(X_q)=c_1(T_{\CP^1})[\CP^1]=\chi(S^2)=2.
\end{equation}
For \(\sltwo\), the canonical adjoint-trace Killing form obeys
\(B(A,A)=4\Tr_{2\times2}(A^2)\), so the displayed matrix gives
\begin{equation}
 B(A_q,A_q)=2\Delta(p).
 \label{eq:canonicalBdelta}
\end{equation}
No rescaling of Lie-algebra generators is allowed here.  Instead, write the
 same matrix in Route E's spherical basis as
\begin{equation}
 A_q=x_+T_{+1}+x_0T_0+x_-T_{-1}
 =\begin{pmatrix}
 x_0/2&-x_+/\sqrt2\\ x_-/\sqrt2&-x_0/2
 \end{pmatrix},\qquad
 a=\frac{x_-}{\sqrt2},\quad b=-x_0,\quad
 c=\frac{x_+}{\sqrt2}.
 \label{eq:sphericalmap}
\end{equation}
Then \(\Delta=x_0^2-2x_+x_-\), while Eq.~\eqref{eq:Ktr} gives
\(B(\bm x,\bm x)=2\sqrt3\,\bm x^T\Ktr\bm x
=2(x_0^2-2x_+x_-)=2\Delta\).  This is an invertible change of coordinates,
not a change of bracket.  It proves the remaining zero-versus-nonzero
equivalence and fixes the normalization.
\end{proof}

The normalization can be checked without words.  Write
\begin{equation}
 p=\frac{x_-}{\sqrt2}u^2-x_0uv+\frac{x_+}{\sqrt2}v^2,
 \qquad \bm x=(x_+,x_0,x_-)^T.
\end{equation}
Then
\begin{align}
 \Delta(p)
 &=x_0^2-2x_+x_-,\\
 \bm x^T\Ktr\bm x
 &=\frac1{\sqrt3}(x_0^2-2x_+x_-),\\
 &\boxed{B(\bm x,\bm x)=2\Delta(p)
 =2\sqrt3\,\bm x^T\Ktr\bm x.}
 \label{eq:discKexact}
\end{align}
The unnormalized second transvectant is
\begin{equation}
 (p,p)_2=p_{uu}p_{vv}-2p_{uv}^2+p_{vv}p_{uu}
 =8ac-2b^2=-2\Delta(p)=-B(\bm x,\bm x),
\end{equation}
so the same invariant is the transvectant singlet, again up to its declared
normalization.

\begin{remark}[What the theorem does not say]
The theorem welds two Route-E structures exactly.  It does not identify the
two zeros with positive/negative energy, two particles, or two families.  It
also shows why the regular-semisimple qualifier matters: a generic nonzero
section has two distinct centers, but a nilpotent section has one center with
multiplicity two.
\end{remark}

\section{A safe two-polarity bridge: Hopf reduction and moment map}

A creative way to retain the source manuscript's two polarities without a
negative-energy instability is to begin with a complex doublet
\(z=(z_0,z_1)\in\mathbb C^2\).  At fixed norm, \(z^\dagger z=1\), the state
space is \(S^3\).  Quotienting the common phase gives
\begin{equation}
 S^3/U(1)\simeq\CP^1.
 \label{eq:Hopfquotient}
\end{equation}
The Hopf map is
\begin{equation}
 \bm n=\bigl(2\Re z_0^*z_1,\;2\Im z_0^*z_1,\;
 |z_0|^2-|z_1|^2\bigr),\qquad \bm n^2=1.
 \label{eq:Hopfmap}
\end{equation}
For the relative circle action, a normalized Fubini--Study moment map is
\begin{equation}
 \mu([z_0:z_1])=
 \frac{|z_0|^2-|z_1|^2}{|z_0|^2+|z_1|^2},
 \qquad -1\le\mu\le1.
 \label{eq:momentmap}
\end{equation}
Its two fixed points satisfy
\begin{equation}
 p_+=[1:0],\quad\mu(p_+)=+1;
 \qquad p_-=[0:1],\quad\mu(p_-)=-1.
\end{equation}
The numerical endpoint values \(\pm1\) use the displayed generator and the
symmetric additive normalization.  Rescaling the generator by \(c\) rescales
them to \(\pm c\).  What is normalization-independent is that the compact
moment image is bounded and that the two fixed points are opposite extrema
(up to reversing orientation).  This suggests the normalized safe dictionary
\begin{equation}
 E_p,E_n\quad\rightsquigarrow\quad\mu=+1,-1
 \quad\text{(charge/orientation polarity),}
 \label{eq:safesign}
\end{equation}
while the physical Hamiltonian remains bounded below.

With
\begin{equation}
 \omega_{\rm FS}=\frac{\ii\,\dd\xi\wedge\dd\bar\xi}
 {(1+|\xi|^2)^2},
 \qquad \frac1{2\pi}\int_{\CP^1}\omega_{\rm FS}=1,
\end{equation}
geometric quantization of Chern number \(k\ge0\) gives
\begin{equation}
 L\simeq\cO(k),\qquad
 \mathcal H_k=H^0(\CP^1,\cO(k))
 \simeq\operatorname{Sym}^k\mathbb C^2,
 \qquad \dim\mathcal H_k=k+1.
 \label{eq:geomquant}
\end{equation}
If, and only if, the physical bubble reduction also satisfies Route E's
self-carried identification \(L=T_{\CP^1}\simeq\cO(2)\), then \(k=2\) and
\(\dim\mathcal H=3\).

\begin{assumption}[Projective-doublet bridge, BA-P]
The energy-information order parameter is a physical complex doublet; its
fixed-charge reduction is \(\CP^1\); its Berry/prequantum bundle equals
\(T_{\CP^1}\); and its chiral zero-mode operator has
\(\ker D^+\simeq H^0(T)\), \(\ker D^-=0\), with no extra charged zero modes.
\end{assumption}

None of the four clauses follows merely from the Hopf identity
\eqref{eq:Hopfquotient}.  This bridge is therefore a concrete research target,
not a derivation of three families.

\section{Conditional Q-ball completion of an energy bubble}

\subsection{Action, charge, and radial equation}

Replace a literal positive/negative-energy bubble by one complex scalar with a
healthy kinetic term:
\begin{equation}
 S_\Phi=\int\dd^4x\sqrt{-g}\left[
 -g^{\mu\nu}\partial_\mu\Phi^*\partial_\nu\Phi-U(|\Phi|)
 \right],
 \label{eq:qballaction}
\end{equation}
with
\begin{equation}
 U(f)=m^2f^2-\lambda f^4+\frac{\eta}{M^2}f^6,
 \qquad m^2,\lambda,\eta,M^2>0.
 \label{eq:qballpotential}
\end{equation}
The global \(U(1)\) current can be normalized so that for
\begin{equation}
 \Phi(t,r)=f(r)e^{\ii\omega t}
\end{equation}
the charge and energy in flat space are
\begin{align}
 Q&=2\omega\int\dd^3x\,f^2,\label{eq:qcharge}\\
 E&=\int\dd^3x\left[\omega^2f^2+(\bm\nabla f)^2+U(f)\right].
 \label{eq:qenergy}
\end{align}
Variation gives the boundary-value problem
\begin{equation}
 f''+\frac2r f'=\left(m^2-\omega^2-2\lambda f^2
 +\frac{3\eta}{M^2}f^4\right)f,
 \qquad f'(0)=0,\quad f(\infty)=0.
 \label{eq:qballode}
\end{equation}

\begin{proposition}[Existence window for the sextic model]
The thin-wall Q-ball window is nonempty when
\begin{equation}
 0<\omega_0^2:=\min_{f>0}\frac{U(f)}{f^2}<m^2,
 \end{equation}
and for \eqref{eq:qballpotential}
\begin{equation}
 f_0^2=\frac{\lambda M^2}{2\eta},\qquad
 \omega_0^2=m^2-\frac{\lambda^2M^2}{4\eta},
 \qquad \omega_0^2<\omega^2<m^2.
 \label{eq:qwindow}
\end{equation}
\end{proposition}

\begin{proof}
The effective ratio is
\begin{equation}
 \frac{U(f)}{f^2}=m^2-\lambda f^2+\frac\eta{M^2}f^4.
\end{equation}
Differentiation with respect to \(f^2\) gives
\(-\lambda+2\eta f^2/M^2=0\), hence the displayed \(f_0\) and \(\omega_0\).
The upper inequality ensures exponential decay at infinity because the
linearized equation has inverse length \(\sqrt{m^2-\omega^2}\); the lower
inequality makes the effective potential \(U_\omega=U-\omega^2f^2\) negative
somewhere, allowing a nontrivial bounce-like radial trajectory.  This is the
standard Q-ball criterion \cite{Coleman1985}; interacting extensions can
alter the stability range and must be retested \cite{Hamada2024}.
\end{proof}

\subsection{Reproducible benchmark and approximation boundary}

For
\begin{equation}
 m=M=1\GeV,\qquad\lambda=\eta=1,\qquad Q=10^6,
\end{equation}
one finds
\begin{equation}
 f_0=\frac1{\sqrt2}\GeV,\qquad
 \omega_0=\frac{\sqrt3}{2}\GeV=0.8660254\GeV.
\end{equation}
At \(\omega=\omega_0\),
\begin{equation}
 U_\omega(f)=f^2(f^2-f_0^2)^2,
\end{equation}
and the wall tension in the present kinetic convention is
\begin{equation}
 \sigma=2\int_0^{f_0}\sqrt{U_\omega(f)}\,\dd f
 =\frac18\GeV^3.
 \label{eq:sigma}
\end{equation}
Using \(Q=2\omega_0f_0^2(4\pi R^3/3)\) gives
\begin{equation}
 R_0=\left(\frac{3Q}{8\pi\omega_0f_0^2}\right)^{1/3}
 =65.0819045\GeV^{-1}=12.8424157\,\mathrm{fm}.
\end{equation}
Adding the leading surface energy yields
\begin{equation}
 \frac EQ\simeq\omega_0+\frac{4\pi R_0^2\sigma}{Q}
 =0.872678754\GeV<m.
 \label{eq:qbenchmark}
\end{equation}
An independent fixed-step-profile variation minimizes
\begin{equation}
 E(R)=\frac{Q^2}{4f_0^2V(R)}+U(f_0)V(R)+4\pi R^2\sigma,
 \qquad V(R)=\frac{4\pi R^3}{3},
 \label{eq:stepenergy}
\end{equation}
and gives
\begin{equation}
 R=64.9712654\GeV^{-1},\quad
 \omega=0.870457184\GeV,\quad
 E/Q=0.872667434\GeV.
\end{equation}
These are thin-wall/step-profile checks, not an exact numerical solution of
\eqref{eq:qballode}.  The exact profile and fluctuation operator remain a
promotion gate.

\subsection{Exact-symmetry phase no-observable theorem}

\begin{theorem}[No phase clock from an exact Q-ball symmetry]
Let the action and a local observable \(\mathcal O\) be invariant under the
global transformation \(\Phi\mapsto e^{\ii\alpha}\Phi\).  On the stationary
Q-ball orbit \(\Phi=f(r)e^{\ii\omega t}\), \(\mathcal O\) is independent of
time.  Therefore the exact Q-ball phase alone cannot make an invariant mass
or visibility periodically appear and disappear.
\end{theorem}

\begin{proof}
For every \(t\), the field is obtained from \(f(r)\) by the symmetry
transformation with \(\alpha=\omega t\).  Thus
\begin{equation}
 \mathcal O[f(r)e^{\ii\omega t}]
 =\mathcal O[e^{\ii\omega t}f(r)]
 =\mathcal O[f(r)].
\end{equation}
In particular, \(|\Phi|^2\), \(T_{\mu\nu}\), and any invariant pole mass are
stationary.  An observable phase requires a second phase reference or an
explicitly symmetry-breaking spurion.
\end{proof}

This creates a useful no-free-lunch constraint.  Exact \(U(1)\) protects
\(Q\) but hides the absolute phase; a fixed relative-phase term can expose the
phase but then generally gives \(\nabla_\mu j^\mu\ne0\), so absolute Q-ball
stability is replaced by a calculable lifetime.

\section{Route-E weight-two phase and a physical relative phase}

Route E reconstructs
\begin{equation}
 M_R=-m_D^T m_\nu^{-1}m_D,\qquad m_D=vY_\nu.
\end{equation}
Under a uniform nonzero complex rescaling \(Y_\nu\mapsto yY_\nu\),
\begin{equation}
 M_R\mapsto y^2M_R.
\end{equation}
If the positive scale is defined by a Takagi singular value,
\(M_*=\sigma_{\max}(M_R)>0\), then
\begin{equation}
 M_*\mapsto|y|^2M_*.
\end{equation}
For the normalized contact projection this gives the exact covariance
\begin{equation}
 \boxed{\zeta\mapsto\frac{y^2}{|y|^2}\zeta
 =e^{2\ii\arg y}\zeta.}
 \label{eq:zetacov}
\end{equation}
Consequently \(\zeta\) is invariant under any nonzero real uniform scaling,
not only positive scaling; for a complex scaling its phase has weight two.
The normalized construction is undefined at \(y=0\), and family-dependent
rescalings do not obey \eqref{eq:zetacov}.

The benchmark is
\begin{equation}
 \zeta=0.1076472949+0.0736514853\ii,\qquad
 |\zeta|=0.1304319033,\qquad \arg\zeta=0.6000380203.
 \label{eq:zetabench}
\end{equation}
Uniform phase covariance is not itself observable.  If \(\Phi\) carries
weight one and \(\zeta\) weight two, a rephasing-invariant relative phase is
\begin{equation}
 e^{\ii\delta}:=
 \frac{\zeta^*\Phi^2}{|\zeta|\,|\Phi|^2},
 \qquad \delta=2\arg\Phi-\arg\zeta.
 \label{eq:relativephase}
\end{equation}
A bounded visibility ansatz with the source manuscript's four special phases
is then
\begin{equation}
 w(\delta)=\frac{1+\cos\delta}{2}
 =\cos^2\!\left(\arg\Phi-\frac12\arg\zeta\right),
 \qquad 0\le w\le1.
 \label{eq:visibility}
\end{equation}
This changes a coupling or spectral residue, not the Poincar\'e Casimir.  The
two naive identifications are numerically inequivalent:
\begin{equation}
 \cos^2(\arg\zeta)=0.68114344,\qquad
 \cos^2(\arg\zeta/2)=0.91265707.
\end{equation}
Route E selects neither convention.

For two complex symmetric Majorana structures
\(M_R=M_V+M_C\), \(M_C=\zeta\Ktr\), a family-basis congruence sends both as
\(M\mapsto U^TMU\).  If \(M_V\) is invertible,
\begin{equation}
 \mathcal R:=M_V^{-1}M_C\mapsto U^{-1}\mathcal R U.
\end{equation}
Therefore quantities such as
\begin{equation}
 \Im\Tr(\mathcal R^n),\qquad \arg\det(I+\mathcal R)
 \label{eq:phaseinvariants}
\end{equation}
are basis invariants that can test a \emph{relative} contact phase.  Any
future phase claim must be written in such invariant form rather than as an
absolute \(\arg\zeta\).

\section{Candidate Q-ball--Schur-complement bridge}

\subsection{Action-level construction}

Let \(N\in V\) and a sterile heavy messenger \(X\in V^*\) be left-handed Weyl
family vectors.  A non-supersymmetric quadratic mass sector with a dynamical
scalar-dependent mixing is
\begin{align}
 \mathcal L_{NX}={}&
 \ii N^\dagger\bar\sigma^\mu D_\mu N
 +\ii X^\dagger\bar\sigma^\mu\partial_\mu X\nonumber\\
 &-\frac{M_*}{2}\left[
 N^TM_VN+2\lambda(\Phi)X^TN-X^T\Ktr^{-1}X
 \right]+\text{h.c.}
 \label{eq:NXaction}
\end{align}
At momenta well below the messenger gap, the algebraic equation is
\begin{equation}
 X=\lambda(\Phi)\Ktr N.
\end{equation}
Substitution gives
\begin{equation}
 \mathcal L_{\rm eff}\supset
 -\frac{M_*}{2}N^T\left[M_V+\lambda(\Phi)^2\Ktr\right]N
 +\text{h.c.}
 \label{eq:schureff}
\end{equation}
More generally, if the messenger kernel is
\(-\mu\Ktr^{-1}\), then
\begin{equation}
 \zeta_{\rm eff}=\frac{\lambda(\Phi)^2}{\mu}.
 \label{eq:zetageneral}
\end{equation}
Route E sets \(\mu=1\) by a field rescaling; keeping it visible is useful for
charge and normalization audits.

The direct phase-doubling bridge is
\begin{equation}
 \lambda(\Phi)=g\frac\Phi\Lambda
 \quad\Longrightarrow\quad
 \boxed{\zeta_{\rm eff}(x)=g^2\frac{\Phi(x)^2}{\Lambda^2}.}
 \label{eq:zetaportal}
\end{equation}
On a Q-ball background,
\begin{equation}
 \zeta_{\rm eff}(r,t)=g^2\frac{f(r)^2}{\Lambda^2}e^{2\ii\omega t}.
 \label{eq:zetaball}
\end{equation}
This exactly reproduces a localized weight-two coefficient.  It is the most
direct candidate relation between the two research programs.

\begin{assumption}[Q-ball contact portal, BA-QK]
The physical heavy mode has the Route-E family representation and kernel in
\eqref{eq:NXaction}; the dynamical mixing is \(g\Phi/\Lambda\); all charge,
gauge, anomaly, and operator-selection rules are satisfied; and the resulting
relative phase cannot be removed by a field redefinition.
\end{assumption}

\subsection{Why the portal is not yet a completion}

The following gates are mandatory.

\paragraph{Charge ledger.}
With \(q_\Phi=+1\), the choices \(q_X=0\), \(q_N=-1\) make \(\Phi XN\) and
\(XX\) invariant, but \(N^TM_VN\) then needs a charge-\(+2\) spurion.  If
\(M_V\) is instead a fixed explicit term, the Q-ball \(U(1)\) is broken.  If
all spurions transform dynamically, the full background may be stationary in
a rotating frame and no visibility modulation follows.  This tension must be
resolved, not hidden in notation.

\paragraph{Tree-level normalization.}
For canonical messenger kinetic terms, the leading low-energy substitution
also gives
\begin{equation}
 Z_N=I+|\lambda|^2\Ktr^\dagger\Ktr
 =\left(1+\frac{|\lambda|^2}{3}\right)I.
 \label{eq:ZN}
\end{equation}
At the Route-E benchmark \(|\lambda|^2=|\zeta|\),
\begin{equation}
 \delta Z_N=\frac{|\zeta|}{3}=0.0434773011.
\end{equation}
This is a tree-level matching effect.  The canonical physical test is the
complete Weinberg combination \(Y M^{-1}Y^T\), not \(\zeta\) by itself.

\paragraph{Full propagator and interaction.}
When external energies are not negligible relative to all heavy Takagi
singular values, one must use
\begin{equation}
 \mathcal M_{NX}=M_*
 \begin{pmatrix}M_V&\lambda I\\ \lambda I&-\Ktr^{-1}\end{pmatrix}
 \label{eq:6x6}
\end{equation}
and its kinetic matrix.  The original fixed-\(\lambda\) model is Gaussian and
does not by itself generate a \(|\lambda|^2/(16\pi^2)\) anomalous dimension.
The dynamical vertex \(g\Phi XN/\Lambda\) supplies a genuine interaction, but
its loop matching, cutoff power counting, and counterterms must be calculated.

\paragraph{Selection and leakage.}
The existing Route-B rule permits the unwanted post-breaking operator
\(XLH\).  A valid discrete/gauge rule must forbid it and enumerate operators
through a declared dimension.  Moreover, if neutrinos or messengers carry the
Q-ball charge, the soliton may emit them.  Absolute stability requires its
energy per unit charge to lie below every kinematically available
charge-carrying channel; explicit phase locking instead requires a lifetime
calculation.

\paragraph{Static benchmark versus rotating background.}
Route E fits a constant \(\zeta\).  Equation \eqref{eq:zetaball} is periodic
relative to a fixed \(M_V\).  It therefore requires a Floquet/adiabatic
seesaw analysis and cannot be inserted into the static fit without rerunning
neutrino, threshold, proton-decay, and cosmology gates.

\section{Candidate dynamical origin of \texorpdfstring{H3$^+$}{H3+}}

The previous sections use H3$^+$ but do not derive it.  There is, however, a
specific loop mechanism worth testing.  Couple a complete heavy family
adjoint multiplet to a background family source \(A_\mu^a\).  Covariance and
the Ward identity constrain its vacuum polarization to
\begin{equation}
 \Pi_{ab}^{\mu\nu}(p)=
 (p^2\eta^{\mu\nu}-p^\mu p^\nu)\,\Pi(p^2)\,
 \Tr_{\rm ad}(\operatorname{ad}T_a\operatorname{ad}T_b).
 \label{eq:adjointloop}
\end{equation}
The family tensor is therefore
\begin{equation}
 B_{ab}^{\rm Kill}:=
 \Tr_{\rm ad}(\operatorname{ad}T_a\operatorname{ad}T_b).
\end{equation}
For a simple algebra, invariant symmetric bilinears form a one-dimensional
space, so a nonzero result is proportional to the Killing form.  On the
Route-E branch this is \(2\sqrt3\Ktr\).

\begin{assumption}[Adjoint-loop H3$^+$ mechanism, BA-H3]
A healthy UV sector contains a complete, non-canceling adjoint current whose
low-energy two-point function induces the required family bilinear, and a
specified source/messenger map converts that bilinear into the Majorana
contact rather than merely a gauge kinetic term.
\end{assumption}

Equation \eqref{eq:adjointloop} explains why a Killing tensor is natural once
an adjoint loop exists; it does not prove the existence of that loop, its
nonzero coefficient, a positive spectral density in the relevant channel, or
the conversion to a Majorana mass.  H3$^+$ therefore remains physically open.

\section{The oscillatory formula: exact solution and strict separation}

The source manuscript includes
\begin{equation}
 \sin(x^2)+\cos(x^2)=\sin(y^2)+\cos(y^2).
 \label{eq:sourceosc}
\end{equation}
Using \(\sin t+\cos t=\sqrt2\sin(t+\pi/4)\), the identity
\(\sin A=\sin B\) gives
\begin{equation}
 A=B+2\pi n\quad\text{or}\quad A=\pi-B+2\pi n,
 \qquad n\in\mathbb Z.
\end{equation}
With \(A=x^2+\pi/4\), \(B=y^2+\pi/4\), the complete real solution set is
\begin{equation}
 \boxed{x^2-y^2=2\pi n}
 \qquad\text{or}\qquad
 \boxed{x^2+y^2=\frac\pi2+2\pi n},
 \label{eq:oscsolutions}
\end{equation}
subject to real-coordinate nonemptiness.  The circular radii are
\begin{equation}
 r_n^2=\frac\pi2+2\pi n,\qquad n\ge0,
\end{equation}
and their exact radial pitch is
\begin{equation}
 r_{n+1}-r_n=\frac{2\pi}{r_{n+1}+r_n}
 \sim\frac\pi{r_n}.
 \label{eq:pitch}
\end{equation}
Thus a shrinking interference pitch is a theorem.  These are continuous
circles/hyperbolae, not discrete atomic orbitals; standard subshell
degeneracies are \(2(2\ell+1)=2,6,10,14,\ldots\).

There is no direct Route-E identification.  The functions
\(\sin(\xi^2)\) and \(\cos(\xi^2)\) have an essential singularity in the
\(\eta=1/\xi\) chart, whereas global sections of \(T_{\CP^1}\) are exactly
degree-at-most-two polynomials.  A shared appearance of squares is
insufficient to identify the state spaces.

\section{Relational time and gravity: parallel interfaces, not Route-E derivations}

The source intuition that time may be relational has a standard conditional
form.  Let a clock \(C\) and system \(S\) satisfy the stationary constraint
\begin{equation}
 (H_C+H_S)|\Psi\rangle=0.
\end{equation}
If clock states obey \(\langle t|H_C=-\ii\hbar\partial_t\langle t|\), then
\begin{equation}
 |\psi(t)\rangle:=\langle t|\Psi\rangle
\end{equation}
satisfies
\begin{align}
 0&=\langle t|(H_C+H_S)|\Psi\rangle\\
 &=-\ii\hbar\partial_t|\psi(t)\rangle+H_S|\psi(t)\rangle,
\end{align}
hence
\begin{equation}
 \ii\hbar\partial_t|\psi(t)\rangle=H_S|\psi(t)\rangle.
 \label{eq:PW}
\end{equation}
This Page--Wootters construction \cite{PageWootters1983} has active modern
extensions, including finite-dimensional time-dilation mechanisms
\cite{Cafasso2024}.  It does not identify time with local \(E_d\), nor does it
derive Route E.

A thermodynamic route to Einstein's equation is likewise conditional.  On
every local Rindler horizon, impose
\begin{equation}
 \delta Q=T_U\delta S,\qquad
 T_U=\frac{\hbar\kappa}{2\pi},\qquad
 \delta S=\frac{\delta A}{4G\hbar}.
\end{equation}
The Raychaudhuri equation relates the area variation to
\(R_{\mu\nu}k^\mu k^\nu\), while the heat flux is proportional to
\(T_{\mu\nu}k^\mu k^\nu\).  Requiring the equality for every null vector
\(k^\mu\) gives
\begin{equation}
 R_{\mu\nu}+\Phi g_{\mu\nu}=8\pi G T_{\mu\nu}.
\end{equation}
Stress conservation and the Bianchi identity fix
\(\Phi=-R/2+\Lambda\), yielding Einstein's equation
\cite{Jacobson1995,Jacobson2016}.  The horizon entropy law and local
equilibrium are inputs.  This route is a gravity interface for later work,
not evidence for the family/contact bridge.

\section{Claim ledger and ordered research program}

\begin{longtable}{>{\raggedright\arraybackslash}p{0.29\textwidth}
                  >{\raggedright\arraybackslash}p{0.18\textwidth}
                  >{\raggedright\arraybackslash}p{0.43\textwidth}}
\toprule
Claim & Status & Boundary / next gate \\
\midrule
\endfirsthead
\toprule
Claim & Status & Boundary / next gate \\
\midrule
\endhead
\(E_d\) alone determines time/gravity & \status{refuted} & Thin-shell and
pressure counterexamples, Eqs.~\eqref{eq:shellmetric}--\eqref{eq:RicciTrace}.\\
Invariant mass disappears at four internal phases & \status{refuted} &
Poincar\'e Casimir; only coupling/residue/gap may be modulated.\\
Route-E CP1/O(2) carrier has three sections & \status{conditional theorem} &
Requires H3$^+$ to select the branch; H3 alone gives only the upper bound.\\
Two centers \(\Leftrightarrow\Delta\ne0\Leftrightarrow\Ktr\)-norm nonzero &
\status{proved} & Exact on the selected CP1 branch,
Theorem~\ref{thm:disc-killing}.\\
Opposite center labels without negative energy & \status{candidate bridge} &
Moment-map polarity is exact; physical identification requires BA-P.\\
Energy bubble as a Q-ball & \status{conditional completion} & Exact radial
solution, Hessian, charge leakage, and gravity remain open.\\
\(\zeta\) carries phase weight two & \status{proved in Route E} & Uniform
complex rescaling only; the absolute phase is not observable.\\
\(\zeta_{\rm eff}=g^2\Phi^2/\Lambda^2\) & \status{candidate bridge} &
Requires BA-QK, charge/selection rules, full matching, and rerun fits.\\
Adjoint loop physically derives H3$^+$ & \status{open mechanism} & Killing
tensor follows conditionally; existence and Majorana conversion do not.\\
Oscillatory equation quantizes atoms or gives CP1 sections & \status{refuted} &
It gives continuous conics and has an essential singularity at CP1 infinity.\\
Relational time / thermodynamic gravity & \status{parallel conditional tracks} &
Useful interfaces; neither derives Route E or the bubble map.\\
\bottomrule
\end{longtable}

The sequential program is deliberately fail-closed.

\begin{enumerate}[label=\textbf{AP--E\arabic*.}]
\item \textbf{Projective-doublet action.}  Write a two-component bubble
action, identify its constraint and gauge/global quotient, and prove or refute
that the physical reduced moduli space is \(\CP^1\).
\item \textbf{Discriminant/contact verifier.}  Preserve
Theorem~\ref{thm:disc-killing} as an exact symbolic regression test, including
the regular/nilpotent boundary and Route-E normalization.
\item \textbf{Self-carried geometric quantization.}  Derive the Berry
curvature and Chern number from AP--E1 rather than imposing them; test whether
the prequantum bundle is actually \(T_{\CP^1}\).
\item \textbf{Chiral zero modes.}  Construct the fluctuation/Dirac operator,
prove \(\ker D^+\simeq H^0(T)\), compute \(\ker D^-\), extra modes, and the
spectral gap.
\item \textbf{Exact Q-ball solution.}  Solve Eq.~\eqref{eq:qballode}, compute
\(Q(\omega)\), \(E(Q)\), fixed-charge Hessian modes, gravitational
compactness, and all charge-emission thresholds.
\item \textbf{Messenger/contact realization.}  Build a gauge- and
anomaly-consistent version of Eq.~\eqref{eq:NXaction}; enumerate operators,
forbid \(XLH\), and verify the full kinetic plus six-by-six matching.
\item \textbf{Physical-phase quotient.}  Remove all field rephasings, list
basis-invariant CP quantities such as Eq.~\eqref{eq:phaseinvariants}, and test
whether a relative phase survives without destroying the Q-ball lifetime.
\item \textbf{H3$^+$ origin test.}  Compute the candidate adjoint current
correlator and its spectral sign; show whether it yields a Majorana contact or
only a kinetic bilinear.
\item \textbf{Integrated phenomenology.}  Rerun Route-E flavor/seesaw,
threshold, proton-decay, cosmology, and time-dependent/Floquet gates with a
single action.  Only this stage can request physics promotion.
\item \textbf{Independent gravity track.}  Constrain any \(E_d\) carrier with
clock redshift, fifth-force, rotation-curve, and lensing data; do not merge it
with AP--E1--9 without an explicit operator connecting the sectors.
\end{enumerate}

\section{2026-07-16 AP--E3 global/nonlinear and AP--E4 SQM handoff}

This section is a reversible index into three standalone derivation notes.
It records the equations needed to reconstruct the decision boundary without
replacing their complete proofs and numerical ledgers.

\subsection{Charged two-colour nonlinear proxy}

The declared six-real-field diagnostic potential is
\begin{equation}
 V=\frac{\lambda}{4}
 \left(\sigma^2+\bm\pi^2+2|\Delta|^2-v^2\right)^2
 -h\sigma+\mu_g^2|\Delta|^2,
 \qquad q_\Delta=2.
 \label{eq:handoff-proxy-potential}
\end{equation}
At the mesonic stationary point, \(h/f=m_\pi^2\), and its exact Hessian
contains
\begin{equation}
 m_\sigma^2=m_\pi^2+2\lambda f^2,\qquad
 m_\Delta^2=m_\pi^2+\mu_g^2.
 \label{eq:handoff-proxy-gaps}
\end{equation}
The requested \((e_g,v_g/\Lambda_c)\) scan is explicitly a matching ansatz,
not a UV loop result:
\begin{equation}
 \frac{\mu_g^2}{\Lambda_c^2}
 =\frac{\mu_{\rm PG,0}^2}{\Lambda_c^2}
 +c_g\left(e_g\frac{v_g}{\Lambda_c}\right)^2.
 \label{eq:handoff-matching-chart}
\end{equation}
The mesonic nonlinear sector solves the \(B=1\) hedgehog with
\begin{equation}
 B=-\frac2\pi\int_0^\infty\sin^2F\,F'\,\dd x=1,
 \qquad E_2-E_4+3E_0=0,
 \label{eq:handoff-b-virial}
\end{equation}
and diagonalizes both the generalized radial operator and the charged block
\begin{equation}
 \cH_{\Delta,\ell}=-\frac{\dd^2}{\dd x^2}+\alpha^2
 +\chi(1-\cos F)+\frac{\ell(\ell+1)}{x^2}.
 \label{eq:handoff-diquark-hessian}
\end{equation}
The deterministic \(18/18\) card passes the classical radial necessary gate.
It deliberately leaves lattice QC2D, the full coupled three-dimensional
Hessian, quantum determinants, and global target-space unwinding open.  The
complete formulas are in
\path{route_f/tex/ap_e3_charged_two_colour_soliton_proxy.tex}.

\subsection{Non-extendible mixed WZW and spin refinement}

Put \(X=S^3\times S^2\), choose integral generators \(u_3,u_2\), and let
\(\widehat\Omega_n\in\widehat H^5(X;\ZZ)\) have curvature
\(n\omega_3\wedge\omega_2\) and characteristic class \(n u_3u_2\).  The
bosonic action on every closed spin four-manifold, without an extension of
its map, is
\begin{equation}
 Z_{\rm bos}(M,\phi)=
 \exp\!\left(2\pi\ii\,\widehat\Omega_n(\phi_*[M])\right).
 \label{eq:handoff-diff-character}
\end{equation}
Here \(H^4(X;\RR/\ZZ)=0\), so fixed integral curvature has no ordinary flat
ambiguity.  Stable splitting instead gives
\begin{equation}
 \Omega_4^{\rm Spin}(S^3\times S^2)
 \simeq\ZZ\oplus\ZZ_2\oplus\ZZ_2.
 \label{eq:handoff-spin-bordism}
\end{equation}
If \(\nu_3\) is the spin-bordism class of the inverse image of a regular
point under the \(S^3\) projection and \(\nu_2\) is the Arf class of the
inverse-image spin surface under the \(S^2\) projection, the general reduced
spin refinement is
\begin{equation}
 Z_{\epsilon_3,\epsilon_2}(M,\phi)=Z_{\rm bos}(M,\phi)
 (-1)^{\epsilon_3\nu_3+\epsilon_2\nu_2}.
 \label{eq:handoff-spin-refinement}
\end{equation}
The bosonic differential-character definition is complete; a microscopic
APS/Dai--Freed determinant must still choose
\((\epsilon_3,\epsilon_2)\) before the spin-refined action is fixed.

\subsection{Pure-\texorpdfstring{\(SU(3)\)}{SU(3)} no-go and near-miss}

Adjoint breaking uses
\begin{equation}
 \iota(A,e^{\ii\alpha})=
 \operatorname{diag}(e^{\ii\alpha}A,e^{-2\ii\alpha}),\qquad
 \ker\iota=\{(1,1),(-1,-1)\},
 \label{eq:handoff-su3-map}
\end{equation}
so every \([SU(2)\times U(1)]/\ZZ_2\) term descending from \(SU(3)\) obeys
\begin{equation}
 2j+q=0\pmod2.
 \label{eq:handoff-parity}
\end{equation}
A colour singlet therefore has even charge, proving that the original
\(\bm1_{+1}\) scalar cannot occur.  Since \(qq(\phi^\dagger)^m\) requires
\(2X_q=mX_\phi\), the desired \(m=2\) would require an odd doublet charge to
equal an even singlet charge.  This is impossible.  The closest explicit
branch instead has
\begin{equation}
 \bm3\to\bm2_{+1}\oplus\bm1_{-2},\qquad
 \overline{\bm3}\to\bm2_{-1}\oplus\bm1_{+2},
 \qquad M_q=m+Vy,\quad M_\ell=m-2Vy.
 \label{eq:handoff-su3-branch}
\end{equation}
Choosing \(m=-Vy\) makes \(M_q=0\) and \(M_\ell=-3Vy\), while an adjoint
scalar mass split can keep \(\bm1_{+2}\) light.  This is an actual tree-level
decoupling but its minimal baryon \(qq\phi^\dagger\) is an \(\cO(1)\)
\(SU(2)_\phi\) flavour doublet, not the Route-E \(\cO(2)\) triplet.  The full
proof is in
\path{route_f/tex/ap_e3_nonextendible_wzw_su3_uv_audit.tex}.

\subsection{Physical tangent fermion and the composability gate}

A half-BPS non-Abelian vortex independently yields a normalizable
\(\CP^1\) orientational modulus and a worldline fermion.  On \(v=1/w\),
\begin{equation}
 \psi^v=-w^{-2}\psi^w,\qquad
 D_t\psi^v=-w^{-2}D_t\psi^w,
 \label{eq:handoff-tangent-fermion}
\end{equation}
so the Grassmann field is physically tangent-valued.  Canonical
Spin\({}^c\) quantization gives
\begin{equation}
 Q_E=\sqrt2\,\bar\partial_E,\qquad
 \cH_{\rm ground}=H^{0,*}(\CP^1,E).
 \label{eq:handoff-sqm}
\end{equation}
In the declared canonical Spin\({}^c\) re-quantization, the tangent Clifford
generator alone corresponds to \(E=\cO\) and gives one ground state.  The
source-selected half-form ordering has no normalizable zero mode.  A
separately derived \(E=\cO(2)\) gives
\begin{equation}
 (h^0,h^1)=(3,0),\qquad
 \lambda_{n,\pm}=\pm\frac2{\sqrt C}\sqrt{n(n+3)},\qquad
 \operatorname{mult}=2n+3,\qquad n\ge1.
 \label{eq:handoff-sqm-spectrum}
\end{equation}
The intrinsic \(27/27\) SQM audit therefore derives a physical tangent
fermion, but does not derive the canonical vacuum line and is not yet in the
charged-two-colour mother model.  AP--E3's signed \(k=+2\) WZW character
composes with it only after proving the same \(B=1\), the same \(\CP^1\)
collective-coordinate map, transgressing the degree-five character along the
spatial three-cycle, and pulling the resulting differential line back with
degree \(+2\).  The
degree-one Route-E portal is postponed until that theorem or an anomaly-free
product compactification is available.  See
\path{route_f/tex/ap_e4_moduli_space_sqm.tex}.

\section{2026-07-16 AP--E completion frontier}

This section integrates the next three requested calculations.  It records
the common normalization and promotion logic; the standalone notes retain
the full proofs, source manifests, and negative controls.

\subsection{\texorpdfstring{Four-dimensional action and graded \(B=1\)
Hessian}{Four-dimensional action and graded B=1 Hessian}}

The declared Euclidean lattice theory has the schematic action
\begin{equation}
 S_E=S_{g,W/{\rm imp}}+S_{\rm gf}+\overline cM_{\rm FP}c
 +S_{\chi\Delta}+\overline\psi D_W[U,V,z]\psi ,
 \label{eq:frontier-lattice-action}
\end{equation}
where \(V\in SU(2)_c\), \(z\in U(1)_g\), the meson \(U\in SU(2)\) is
neutral, and the complex diquark has charge two.  At
\((\overline V,\overline z,\overline U,\overline\Delta,\overline\psi)
=(1,1,U_{B=1},0,0)\), write the even variables as
\(\Phi=(A,b,\eta,d_1,d_2)\).  The complete formal quadratic structure is
\begin{equation}
 \mathbb H_{\rm cl}=
 \begin{pmatrix}
 K_{\Phi\Phi}&0&0\\
 0&0&D_W\\
 0&-D_W^T&0
 \end{pmatrix},
 \qquad
 K_{\Phi\Phi}=
 \begin{pmatrix}
 K_{AA}&K_{Ab}&K_{A\eta}&K_{Ad_1}&K_{Ad_2}\\
 K_{bA}&K_{bb}&K_{b\eta}&K_{bd_1}&K_{bd_2}\\
 K_{\eta A}&K_{\eta b}&K_\eta&K_{\eta d_1}&K_{\eta d_2}\\
 K_{d_1A}&K_{d_1b}&K_{d_1\eta}&K_\Delta&0\\
 K_{d_2A}&K_{d_2b}&K_{d_2\eta}&0&K_\Delta
 \end{pmatrix},
 \label{eq:frontier-super-hessian}
\end{equation}
together with the ghost block.  The one-loop functional and its fermionic
collective-coordinate curvature are
\begin{align}
 \Gamma^{(1)}
 &=\frac12\log\det{}'K_{\Phi\Phi}
   -\log\det{}'M_{\rm FP}-\Re\log\det D_W,
 \label{eq:frontier-one-loop}\\
 \partial_i\partial_j\Gamma_f^{(1)}
 &=\Re\Tr\!\left(
 D^{-1}D_iD^{-1}D_j-D^{-1}D_{ij}\right).
 \label{eq:frontier-logdet-curvature}
\end{align}
This is the requested full \emph{formal} gauge--meson--ghost--fermion
Hessian.  The numerical certificate is deliberately smaller: it uses a
\(2^4\) trivial-link background, substitutes a frozen positive
Gauss--Newton meson block for the full Skyrme Jacobi operator, and evaluates
one fermion curvature direction.  It therefore cannot be called dynamical
lattice QC2D or the complete nonlinear Hessian.

The \(24/24\) card gives
\begin{align}
 B(0)&=0.999981468235,\nonumber\\
 \lambda_{\min}(K_{\eta,{\rm test}})&=0.621297848377,\qquad
 \lambda_{\min}(K_\Delta)=0.694069172619,\nonumber\\
 \lambda_{\min}(K_{\Delta,{\rm control}})&=-0.349308273811,\qquad
 \sigma_{\min}(D_W)=0.702570765924,
 \label{eq:frontier-lattice-numbers}
\end{align}
with gauge/ghost \(\xi\)-corrected residual \(1.42\times10^{-14}\), Schur
solve residual \(6.30\times10^{-16}\), log-determinant curvature relative
error \(1.32\times10^{-8}\), and free Wilson/Symanzik observed orders
\(1.99861/3.99628\).  Monte Carlo, measure positivity, the exact Skyrme
Jacobi zero-mode projection, continuum renormalization, FR quantization, and
finite-amplitude stability all remain false gates.  Full details are in
\path{route_f/tex/ap_e5_4d_qc2d_lattice_hessian.tex}.

\subsection{Two torsion generators and regulator-dependent phases}

Explicit representatives of the reduced bordism generators are
\begin{align}
 G_3&=(S^1_{\rm R}\times S^3,\operatorname{pr}_{S^3},{\rm pt}),
 \nonumber\\
 G_2&=(T^2_{\rm RR}\times S^2,{\rm pt},\operatorname{pr}_{S^2}).
 \label{eq:frontier-bordism-generators}
\end{align}
Their transverse inverse images are the nonbounding periodic circle and the
odd-spin torus.  Hence
\begin{equation}
 \ind_2D_{S^1_{\rm R}}=1,\qquad
 \ind_2D^+_{T^2_{\rm RR}}=\Arf(T^2_{\rm RR})=1.
 \label{eq:frontier-mod-two}
\end{equation}
The chirality-paired ambient four-dimensional product spectra define only a
factorized reference phase \((+1,+1)\).  Stacking the two explicitly
localized defect regulators instead gives the complete character table
\begin{equation}
 Z_{\epsilon_3,\epsilon_2}(G_3)=(-1)^{\epsilon_3},
 \qquad
 Z_{\epsilon_3,\epsilon_2}(G_2)=(-1)^{\epsilon_2}.
 \label{eq:frontier-character-table}
\end{equation}
Thus all four torsion assignments are realizable with the same de Rham
curvature.  This proves regulator non-uniqueness; it does not select the
charged-QC2D phase.  That selection still requires the complete
Dirac--Yukawa operator, heavy mass signs, and its odd-dimensional
mapping-torus APS/Dai--Freed problem.

\subsection{A semisimple replacement for the diagonal quotient}

For
\([SU(2)_c\times U(1)]/\mathbb Z_2\), every representation obeys
\begin{equation}
 2j+q=0\pmod2.
 \label{eq:frontier-global-parity}
\end{equation}
It admits \(\bm2_{+1}\) but forbids \(\bm1_{+1}\); hence any successful
candidate must retain the unquotiented global form.  Besides the minimal
backup \(SU(2)_c\times SU(2)_H\), the scan finds the simply connected simple
group
\begin{equation}
 Sp(4)\equiv USp(4)\simeq Spin(5)
 \label{eq:frontier-sp4}
\end{equation}
with branching
\begin{align}
 \bm4&\longrightarrow
 \bm2_0\oplus\bm1_{+1}\oplus\bm1_{-1},\nonumber\\
 \bm5&\longrightarrow
 \bm2_{+1}\oplus\bm2_{-1}\oplus\bm1_0,\nonumber\\
 \bm{10}&\longrightarrow
 \bm1_{+2}\oplus\bm2_{+1}\oplus\bm3_0\oplus\bm1_0
 \oplus\bm2_{-1}\oplus\bm1_{-2}.
 \label{eq:frontier-sp4-branching}
\end{align}
Two Dirac \(\bm5\)'s provide the two light quark flavours.  Two, not one,
complex scalar \(\bm4\)'s are required so their copy index retains the
global \(SU(2)_\phi\) doublet.  Then
\begin{equation}
 qq(\phi_I^\dagger\phi_J^\dagger)_{\rm sym}
 \in\bm3_\phi
 \label{eq:frontier-sp4-operator}
\end{equation}
is the charge-neutral exactly-two \(\mathcal O(2)\) operator at the
representation-theory level.  This does not rederive the earlier unit-cell
Gauss rule.

One real adjoint \(\bm{10}\) breaks the second \(SU(2)\) to \(U(1)\).
For \(M_q=m+q\,yV\), choosing \(m=-yV\) keeps \(q=+1\) light and makes the
\(q=-1,0\) partners heavy at tree level.  The displayed field content has
\begin{equation}
 b_0^{Sp(4)}
 =11-\frac43(2)-\frac16(3)-\frac13(1)
 =\frac{15}{2}>0.
 \label{eq:frontier-sp4-beta}
\end{equation}
Vectorlike and Witten-parity checks pass.  The tuning is not radiatively
protected, however, and neither the full
\(\Omega_5^{\rm Spin}(BSp(4))\) character nor the heavy determinant
threshold has been computed.  The algebraic light coefficient
\(N_cX_q=+2\) is therefore not yet an all-scale orientation theorem.  See
\path{route_f/tex/ap_e3_aps_global_form_search.tex}; its card passes
\(58/58\).

\subsection{Same-soliton SQM, Callias line, CPT, and WZW descent}

For a hypothetical same-soliton Fredholm family
\begin{equation}
 \mathscr D_m^+=D_{A(m)}^++\mathscr Y_m,
 \qquad m\in\mathcal M_1\simeq\mathbb{CP}^1,
 \label{eq:frontier-callias-family}
\end{equation}
let the positive-mass eigenline on
\(S^2_\infty\times\mathbb{CP}^1\) satisfy
\begin{equation}
 c_1(\mathcal E_+)=p x+q y.
 \label{eq:frontier-positive-line}
\end{equation}
The families-Callias push-forward gives, in the fixed orientation
convention,
\begin{equation}
 \operatorname{rank}\Ind\mathscr D^+=\epsilon p,\qquad
 c_1(\det\Ind\mathscr D^+)=\epsilon p q\,y.
 \label{eq:frontier-callias-pushforward}
\end{equation}
Thus a rank-one \(\mathcal O(+2)\) determinant template requires
\(\epsilon p=1,q=2\).  The finite Berry/kernel-family regression obtains
\(c_1=1.999999999999981\) with unit nonzero gap.  It is a topology
template, not the missing four-dimensional Yukawa operator or its Callias
index.

In one fixed Dolbeault polarization, the conjugate triplet is fixed by
Serre duality:
\begin{equation}
 E_-=
 K_{\mathbb{CP}^1}\otimes E_+^\vee
 =\mathcal O(-2)\otimes\mathcal O(-2)
 =\mathcal O(-4),
 \qquad (h^0,h^1)=(0,3).
 \label{eq:frontier-cpt}
\end{equation}
The naive inverse \(\mathcal O(-2)\) instead has \((h^0,h^1)=(0,1)\).
Equation~\eqref{eq:frontier-cpt} closes the mathematical fixed-polarization
map, not its microscopic CPT regulator.

On framed configuration space, spatial fiber integration is intrinsically
defined:
\begin{equation}
 \widehat{\mathcal L}_B
 =\int_{\Sigma_3}\operatorname{ev}^*\widehat\Omega_5
 \in\widehat H^2(\widetilde{\mathcal C}_B;\mathbb Z).
 \label{eq:frontier-transgression}
\end{equation}
It descends through
\(p:\widetilde{\mathcal C}_B\to\mathcal C_B\) only when
\begin{equation}
 \widehat{\mathcal L}_B\in\operatorname{im}p^*,
 \label{eq:frontier-descent}
\end{equation}
which requires an equivariant differential refinement, horizontal invariant
curvature, trivial vertical and large-gauge holonomy, and trivial stabilizer
characters.  After that descent, its pullback to a same-soliton
\(\mathbb{CP}^1\) is \(\mathcal O(+2)\) precisely when
\begin{equation}
 \int_{\Sigma_3\times\mathbb{CP}^1}
 \operatorname{ev}^*c(\widehat\Omega_5)=+2.
 \label{eq:frontier-transgression-number}
\end{equation}
The factorized conditional map yields \(nB\deg(s)=2\), but the present
theory has not identified its scalar \(S^2\) with a normalizable soliton
modulus and has not proved Eq.~\eqref{eq:frontier-descent}.  The same-model
card therefore passes \(27/27\) mathematical controls while keeping every
composition flag false.  The full derivation is
\path{route_f/tex/ap_e4_same_soliton_callias_descent.tex}.

\subsection{Portal authorization theorem and current decision}

Let
\begin{align}
 C_{\rm lat}&=\text{dynamical 4d phase plus complete projected quantum
 Hessian},\nonumber\\
 C_{\rm APS}&=\text{selected torsion phases plus protected semisimple
 all-scale completion},\nonumber\\
 C_{\rm SQM}&=\text{same-soliton SQM, Callias/CPT line, and gauge-basic
 WZW descent}.
 \label{eq:frontier-route-booleans}
\end{align}
The user-specified ordering is encoded exactly as
\begin{equation}
 {\rm portal\ authorized}
 \quad\Longleftrightarrow\quad
 C_{\rm lat}\lor C_{\rm APS}\lor C_{\rm SQM}.
 \label{eq:frontier-portal-authorization}
\end{equation}
All three booleans are currently false.  Consequently no portal operator is
written.

When Eq.~\eqref{eq:frontier-portal-authorization} first becomes true, the
portal theorem must still exhibit a map
\begin{equation}
 F:\mathcal M_1\longrightarrow\Sigma_E,\qquad
 \deg F=+1,\qquad
 F^*\mathcal O_{\Sigma_E}(2)\simeq\mathcal O_{\mathcal M_1}(2),
 \label{eq:frontier-degree-one}
\end{equation}
fix its orientation, show that its operator is gauge and anomaly consistent,
and keep every coupling below the retained bosonic, fermionic, and Landau
gaps.  It must also preserve the physical CPT pairing and introduce no
extra zero modes.  These are recorded as future prerequisites, not assumed
facts.

The integrated promotion card
\path{route_f/code/verify_ap_e5_completion_frontier.py} passes \(20/20\) and
sets
\begin{equation}
 \boxed{\begin{gathered}
 \texttt{any\_preportal\_route\_closed=false},\\
 \texttt{degree\_one\_portal.constructed=false},\\
 \texttt{physics\_promotion\_allowed=false}
 \end{gathered}}
 \label{eq:frontier-final-status}
\end{equation}

\section{AP-E6 same-background audit: exact obstructions and open gates}
\label{sec:ap-e6-same-background}

This section supersedes the completion-frontier status immediately above.
It records the next execution round without weakening the portal ordering.
The compact group denoted \(Sp(4)\) in the physics convention is
\(USp(4)\), or \(Sp(2)\) in the mathematical convention used in classifying
spaces.

\subsection{The exact Sp(4) gauge-bordism result}

The low-degree integral cohomology and spin-bordism coefficients are
\begin{align}
 H^*(BSp(2);\ZZ)&=\ZZ[q_1,q_2],& |q_1|&=4,&|q_2|&=8,\nonumber\\
 (\Omega_q^{\rm Spin})_{q=0}^5
  &=(\ZZ,\ZZ_2,\ZZ_2,0,\ZZ,0).
 \label{eq:e6-bsp-input}
\end{align}
In the homological Atiyah--Hirzebruch spectral sequence
\begin{equation}
 E^2_{p,q}=H_p(BSp(2);\Omega_q^{\rm Spin})
 \Longrightarrow \Omega_{p+q}^{\rm Spin}(BSp(2)),
 \label{eq:e6-bsp-ahss}
\end{equation}
the only nonzero term with \(p+q=5\) is
\begin{equation}
 E^2_{4,1}=H_4(BSp(2);\ZZ_2)=\ZZ_2.
\end{equation}
The possible \(d_2\) source and target use \(H_6\) and \(H_2\), both zero;
the \(d_3\) target uses \(H_1\), and the only later possible map from this
torsion group to \(E^r_{0,4}\simeq\ZZ\) is zero.  There is no extension
ambiguity on this diagonal.  Hence
\begin{equation}
 \boxed{\Omega_5^{\rm Spin}(BSp(2))\simeq\ZZ_2.}
 \label{eq:e6-bsp-result}
\end{equation}
A geometric representative is a unit \(Sp(2)\) instanton on \(S^4\) times
the nonbounding spin circle.  Restriction to the instanton \(SU(2)_c\)
gives
\begin{align}
 \bm4&\downarrow \bm2\oplus\bm1\oplus\bm1,&
 \bm5&\downarrow \bm2\oplus\bm2\oplus\bm1,\nonumber\\
 \bm{10}&\downarrow \bm3\oplus\bm2\oplus\bm2
                   \oplus\bm1\oplus\bm1\oplus\bm1.
\end{align}
With the doublet index normalized to one, the instanton indices are
\((1,2,6)\).  The corresponding mod-two characters therefore are
\begin{equation}
 (\chi_{\bm4},\chi_{\bm5},\chi_{\bm{10}})=(-1,+1,+1).
 \label{eq:e6-sp4-characters}
\end{equation}
Thus the displayed two Dirac \(\bm5\)'s are trivial on the unique gauge
bordism generator.  This closes the gauge-bordism subgate only.

It does not select the two target-space signs.  Independently,
\begin{equation}
 \widetilde\Omega_4^{\rm Spin}(S^3\times S^2)
 \simeq \ZZ_2\{S^1_R\times S^3\}
 \oplus\ZZ_2\{T^2_{\Arf=1}\times S^2\}.
 \label{eq:e6-target-bordism}
\end{equation}
Without a common microscopic mass-family lift and APS boundary problem,
invertible torsion sectors can flip either character.  The gauge character
in Eq.~\eqref{eq:e6-sp4-characters} and the ordered target pair are different
questions.

The Euclidean note verifies gamma-five reality, two-flavour determinant
positivity away from zeroes, and the Pauli--Villars moment identities
\begin{equation}
 c_a=(1,-3,3,-1),\qquad
 \sum_a c_a a^j=0\quad(j=0,1,2).
 \label{eq:e6-pv-moments}
\end{equation}
It does not specify the full interpolation \(K(t)\), regulator statistics,
APS domain, finite local counterterms, or regulated gauge--scalar--ghost
complex.  Accordingly these are partial controls, not a complete Euclidean
regulator.

At the chosen light branch the infrared orientation remains
\begin{equation}
 k_{\rm IR}=N_cX_qB=2\cdot1\cdot1=+2.
 \label{eq:e6-ir-orientation}
\end{equation}
For a uniform flavour action on a complete \(\bm5\), however, the algebraic
charge trace is only the conditional obstruction
\begin{equation}
 k_{\rm full}=2(+1)+2(-1)+0=+2-2=0.
 \label{eq:e6-threshold-obstruction}
\end{equation}
The local pure-gauge shifts are periodic, but no unbroken-group or
gravitational APS determinant ratio has been evaluated.  Therefore
Eq.~\eqref{eq:e6-threshold-obstruction} is not an actual heavy-determinant
matching computation and neither target eta sign is selected.

\subsection{One genuinely relaxed continuum B=1 field, but an off-shell
cubic Hessian}

The declared classical sector is enlarged to a unit \(O(4)\) meson field
\(n\), a neutral breathing field \(s\), and a quadratic charge-two field.
Its continuum energy is
\begin{align}
 E[n,s]=\int\!\dd^3x\bigg\{&\frac12|\partial_i n|^2
 +\frac\kappa2\sum_{i<j}
 \bigl(|\partial_i n|^2|\partial_jn|^2
 -(\partial_i n\!\cdot\!\partial_jn)^2\bigr)\nonumber\\
 &+\frac12|\nabla s|^2+\frac12m_s^2(s-1)^2
 +\beta^2s^2(1-n_0)\bigg\}.
 \label{eq:e6-coupled-energy}
\end{align}
For \(n=(\cos F,\widehat{\bm x}\sin F)\), direct variation gives
\begin{align}
 (r^2+2\kappa\sin^2F)F''={}&-2rF'
 -2\sin F\cos F(\kappa F'^2-1)
 +\frac{2\kappa\sin^3F\cos F}{r^2}
 +\beta^2r^2s^2\sin F,\nonumber\\
 s''+\frac2r s'={}&m_s^2(s-1)+2\beta^2s(1-\cos F).
 \label{eq:e6-coupled-bvp}
\end{align}
With \((\beta,m_s,\kappa)=(1/2,5/2,1)\), a nonlinear BVP solve on
\(R=8,10,12,16\) converges to
\begin{equation}
 E/(4\pi)=6.1282155,\qquad s(0)=0.91712896,\qquad B=+1,
 \label{eq:e6-bvp-anchor}
\end{equation}
and obeys the generalized Derrick identity.  The canonical 4097-point
little-endian float64 array \((r,F,s)\) has SHA-256
\begin{center}
\texttt{81104f59a5f3f4337739fc2a217cafc8da91581c9f1fb40b4fdba6ce0f948d59}.
\end{center}
This proves stationarity within the coupled hedgehog sector.  It is not a
Palais symmetric-criticality proof for every three-dimensional variation.

For the separately declared cubic action, the analytic constrained Hessian
is
\begin{equation}
 H_{\rm tan}=T^T\left[H_{\rm emb}
 -\bigoplus_x(n_x\!\cdot g_x)\bm1_4\right]T,
 \qquad
 H_{\rm proj}=P_{\rm coll}H_{\rm tan}P_{\rm coll}.
 \label{eq:e6-projected-hessian}
\end{equation}
Every \(n+s\) entry is differentiated from one finite-difference energy;
sparse symmetry, geodesic second differences, and the
translation/isorotation projector are checked independently.  The quadratic
diquark and free toron/ghost operators use separately declared grids and
boundaries.  They are audited controls, not blocks of one aggregate matrix.
Thus the \(n+s\) second variation is exact while the complete declared-sector
sparse-Hessian gate remains false.

It is not yet a physical stability operator.  At fixed \(L=8\), the sampled
continuum solution gives
\begin{center}
\begin{tabular}{rrrr}
\toprule
\(N\)&\(a\)&\(B_a\)&
\(\|\operatorname{grad}E_a\|_{2,{\rm dens}}\)\\
\midrule
9&1.000000&0.168018&0.34766\\
13&0.666667&0.440012&0.43018\\
17&0.500000&0.625354&0.42171\\
\bottomrule
\end{tabular}
\end{center}
The background is therefore off shell for the cubic action and its topology
estimator has not converged to one.  The observed negative projected
curvatures are useful falsifiers of premature stability claims, but they
cannot be promoted to soliton instability eigenvalues.  No interacting
colour links, fermion determinant, four-dimensional importance sampling, or
continuum quantum limit is present.
Moreover, the cubic boundary samples the finite \(R=16\) profile rather than
imposing the exact vacuum; at \(L=8\) its maximum boundary deviation is about
\(4.93\times10^{-2}\).  This boundary mismatch is recorded as part of the
off-shell diagnostic.

\subsection{The same canonical field and the physical-moduli obstruction}

The Yukawa/Callias audit consumes the checksum above rather than re-solving
an older profile.  The declared constituent Yukawa term couples to \(F\) but
not to \(s\); this decoupling is part of the action, not a silent profile
replacement:
\begin{equation}
 \mathcal D_E=\gamma^\mu D_\mu
 +M(P_RU_F+P_LU_F^\dagger),\qquad
 H_F=-\ii\alpha^iD_i+\beta M
 \bigl(\cos F+\ii\gamma_5\bm\tau\!\cdot\!\widehat{\bm x}\sin F\bigr).
 \label{eq:e6-yukawa-operator}
\end{equation}

The true localized orbit is found before any family index is assigned.  If
\(A\in SU(2)_V\) fixes \(U_F\), then at a radius with \(\sin F\ne0\),
\(A(\bm\tau\cdot\widehat{\bm x})A^{-1}
=\bm\tau\cdot\widehat{\bm x}\) for every direction.  Hence \(A\) commutes
with all three Pauli matrices and
\begin{equation}
 \operatorname{Stab}_{SU(2)_V}(U_F)=\ZZ_2,\qquad
 \mathcal M_{\rm rot}=SU(2)/\ZZ_2\simeq SO(3),
 \label{eq:e6-actual-moduli}
\end{equation}
not \(\CP^1\).  Quotienting combined space--isospin rotations by the
diagonal hedgehog stabilizer gives the same \(SO(3)\), not a second orbit.
The separate scalar-vacuum \(\CP^1\) is spatially uniform, with
\begin{equation}
 \|\delta_m\phi\|^2_{[-L,L]^3}=8L^3|\delta_m\phi|^2,
 \label{eq:e6-bulk-cp1-norm}
\end{equation}
so it is a bulk Goldstone direction rather than a normalizable soliton
modulus.

At infinity Eq.~\eqref{eq:e6-yukawa-operator} has a constant invertible mass.
Its positive boundary eigenbundle on \(S^2_\infty\) is therefore trivial,
\begin{equation}
 c_1(E_+)=0,\qquad \ind_{\rm Callias}=0.
 \label{eq:e6-callias-zero}
\end{equation}
This static statement does not contradict spectral flow along a
four-dimensional interpolation; spectral flow equal to \(B\) does not force
an endpoint zero mode.

There is also a representation-level obstruction to the former two-band
template.  A globally nondegenerate \(2\times2\) Hermitian mass with one
positive band on \(X=S^2_\infty\times\CP^1\) defines
\(f:X\to\CP^1\), so its eigenline satisfies
\begin{equation}
 c_1(E_+)^2=f^*(u^2)=0.
\end{equation}
The desired class instead obeys
\begin{equation}
 c_1(E_+)=x+2y,\qquad c_1(E_+)^2=4xy\ne0.
 \label{eq:e6-two-band-obstruction}
\end{equation}
Thus it requires at least three trivial bands, a nontrivial ambient bundle,
or a genuinely infinite-dimensional family, followed by a new uniform-gap
proof.  On the actual \(SO(3)\), \(H^2(SO(3);\ZZ)=\ZZ_2\), so only a torsion
line, not a free \(\cO(2)\) class, is available.

Finally, the correctly typed universal target character is
\begin{equation}
 \widehat\Xi_5
 =2\,\operatorname{pr}_U^*\widehat\omega_3
 \smile\operatorname{pr}_s^*\widehat c_1(\mathsf H)
 \in\widehat H^5(SU(2)\times\CP^1;\ZZ).
 \label{eq:e6-universal-character}
\end{equation}
Only after pulling it back by the evaluation map may one form
\begin{equation}
 \int_{\Sigma_3}\operatorname{ev}^*\widehat\Xi_5
 \in\widehat H^2(\mathcal C_{\rm restricted};\ZZ).
 \label{eq:e6-typed-transgression}
\end{equation}
This establishes a typed local ansatz.  It does not yet construct the full
equivariant differential cocycle, prove all vertical and large-gauge
holonomies, extend through Higgs zeroes and defects, or pull an \(\cO(2)\)
line back to a physical same-soliton \(\CP^1\).  The finite CPT determinant
identity is likewise an algebraic conjugation control, not a microscopic
global regulator.

\subsection{AP-E6 decision}

With the stronger definitions forced by this round, let
\begin{align*}
 C_{Sp}&=(\hbox{complete regulator})\wedge(\hbox{both target eta signs})\\
 &\quad\wedge(\hbox{full heavy APS match})\wedge(\hbox{strong phase}),\\
 C_{B1}&=(\hbox{same stationary lattice solution})\wedge(B_a\to1)\\
 &\quad\wedge(\hbox{full projected super-Hessian})
 \wedge(\hbox{4d dynamics}),\\
 C_{\rm Cal}&=(\hbox{physical moduli})\wedge(\hbox{actual gapped family})\\
 &\quad\wedge(\hbox{CPT determinant line})
 \wedge(\hbox{gauge-basic descent}).
\end{align*}
The exact gauge-bordism subgate, continuum BVP, sparse differentiation, and
moduli/two-band no-go are genuine advances, but
\begin{equation}
 C_{Sp}=C_{B1}=C_{\rm Cal}=\mathrm{false},\qquad
 \boxed{\text{portal not started}.}
 \label{eq:e6-final-decision}
\end{equation}
The canonical standalone records are
\path{route_f/tex/ap_e6_sp4_eta_bordism_threshold.tex},
\path{route_f/tex/ap_e6_relaxed_b1_multigrid_hessian.tex}, and
\path{route_f/tex/ap_e6_same_soliton_yukawa_callias.tex}, together with their
machine-readable cards and the independent AP-E6 frontier verifier.

\section{AP-E7 regulator, discrete topology, and determinant-family audit}
\label{sec:ap-e7-audit}

This section supersedes the AP-E6 execution status while retaining every
negative boundary proved there.  It executes the three requested alternatives
in order.  The result is deliberately asymmetric: the common quadratic
regulator, a pure heavy threshold, the finite-site topology theorem, the
classification of the physical \(SO(3)\) collective-coordinate lines, and a
minimal rank-three mass family are exact.  None of them by itself closes a
same-model pre-portal physics lane.

\subsection{Common quadratic APS/PV complex and its exact scope}

For a smooth background \(\bar X=(\bar A,\bar\Sigma,\bar\Phi)\), let the
infinitesimal gauge action be
\begin{equation}
 {\cal R}_{\bar X}\epsilon
 =\left(D_{\bar A}\epsilon,[\bar\Sigma,\epsilon],-\epsilon\bar\Phi\right).
 \label{eq:e7-gauge-action}
\end{equation}
Using background Feynman gauge, the complete bare quadratic boson/ghost
operators are defined without deleting mixing terms:
\begin{equation}
 \Delta_B=S_{\rm bos}^{(2)}[\bar X]
       +{\cal R}_{\bar X}{\cal R}_{\bar X}^{\dagger},\qquad
 \Delta_{\rm gh}={\cal R}_{\bar X}^{\dagger}{\cal R}_{\bar X}.
 \label{eq:e7-quadratic-complex}
\end{equation}
They act on the stated \(H^2\) Sobolev domains after background stabilizers
are projected.  The fermion amplitude is built from
\(H=\gamma_5{\cal K}\).  Spinor charge conjugation is retained only as a
reference/scalar-mass algebraic control; no fixed-Higgs Kramers symmetry is
assumed for a general adjoint mass.

The common determinant prescription uses
\begin{equation}
 c_a=(1,-3,3,-1),\qquad \Lambda_a^2=a\Lambda^2,\qquad a=0,1,2,3,
 \label{eq:e7-pv-data}
\end{equation}
for which
\begin{equation}
 \sum_ac_aa^r=0\quad(r=0,1,2),\qquad
 \sum_ac_aa^3=-6,\qquad
 \sum_ac_a\log(\lambda+a\Lambda^2)
 =-\frac{2\Lambda^6}{\lambda^3}+O(\lambda^{-4}).
 \label{eq:e7-pv-moments}
\end{equation}
At the level of zeta determinants the regulated quadratic amplitude is
\begin{align}
 Z_{\rm quad}^{\rm PV}(\Lambda)={}&\tau_F
 \prod_{a=0}^3
 \det{}'_{\theta_B}(\Delta_B+a\Lambda^2)^{-c_a/2}
 \det{}'(\Delta_{\rm gh}+a\Lambda^2)^{c_a}\nonumber\\
 &\times\det{}'(H^2+a\Lambda^2)^{N_fc_a/2},\qquad N_f=2.
 \label{eq:e7-common-pv}
\end{align}
Here \(\theta_B=\pi-\varepsilon\) is the upper lateral cut at a saddle.
Equation~\eqref{eq:e7-common-pv} is a determinant prescription; it is not a
claim that every half-integral power has already been realized by a local
second-order Grassmann PV field.  The finite part subtracts the integrated
supercomplex \(A_4\) coefficient at a fixed scale and declares positive
reference masses to have zero local gauge, Hopf, and gravitational theta
angles.

On \(Y_5=M_4\times[0,1]\), the interpolation operator and its literal APS
domain are
\begin{align}
 {\cal D}_5^+&=\partial_t+H(t),\nonumber\\
 \operatorname{Dom}_{\rm APS,0}({\cal D}_5^+)
 =\{u\in H^1(Y_5):{}&P_{[0,\infty)}(H_0)u|_0=0,\quad
 P_{(-\infty,0)}(H_1)u|_1=0\}.
 \label{eq:e7-aps-domain}
\end{align}
The physical determinant-sign formula uses the fixed cut \(0\) and gapped
endpoints:
\begin{equation}
 \tau_F[X(t)]=\exp\!\left(\ii\pi N_f
        \operatorname{SF}_0\{H(t)\}\right).
 \label{eq:e7-physical-sf}
\end{equation}
An arbitrary common resolvent cut still defines a Fredholm APS problem, but
its spectral flow cannot replace Eq.~\eqref{eq:e7-physical-sf} without the
corresponding endpoint spectral-section orientation term.  This distinction
is part of the regulator, not a convention that may be discarded.

The constant negative-mass heavy bundle is
\begin{equation}
 E_h=(V_c\otimes L^{-1})\oplus\bm1,\qquad x=c_1(L),
\end{equation}
so the one-flavour index is
\begin{equation}
 I_h=\int_{M_4}\widehat A(TM_4)\operatorname{ch}(E_h)\big|_4
 =\int_{M_4}\left[-c_2(V_c)+x^2-\frac18p_1(TM_4)\right]\in\ZZ.
 \label{eq:e7-heavy-index}
\end{equation}
The rank-two charged branch contributes \(-p_1/12\), and the neutral singlet
contributes \(-p_1/24\); hence gravity has not been omitted.  For two identical
Dirac flavours,
\begin{equation}
 \boxed{\tau_h^{\rm pure}=\exp(\ii\pi\,2I_h)=1}
 \label{eq:e7-heavy-phase}
\end{equation}
on every closed spin four-manifold and every unbroken
\(SU(2)_c\times U(1)_H\) bundle.  This closes the pure unbroken-gauge/gravity
threshold, not a target-dependent mixed WZW determinant.

On the separated Higgs orbit, one can consistently declare
\begin{equation}
 h=\widehat\Sigma_{\bm5}/V,\qquad
 P_\pm=\frac{h^2\pm h}{2},\qquad P_0=1-h^2,\qquad
 (\rank P_+,\rank P_-,\rank P_0)=(2,2,1).
 \label{eq:e7-higgs-projectors}
\end{equation}
A GW-compatible target source built from these projectors has even copies on
both target generators and therefore conditionally yields \((+1,+1)\) in the
no-stacked-sector convention.  It is not promoted to a finite overlap
Dirac--Yukawa theorem until the ordinary/GW chiral projectors, five-dimensional
kernel, and determinant orientation are specified in one operator.

There is a sharp reason this source does not restore the original deep-UV
claim.  Suppose an exact separated rank-two projector \(C(\Sigma)\) were
continuous and \(Sp(4)\)-equivariant through \(\Sigma=0\).  Since the complex
\(\bm5\) is irreducible, Schur's lemma gives
\begin{equation}
 C(0)=c\bm1_5,\qquad \Tr(HC(0))=0,\qquad
 \Tr(HP_+)=2.
 \label{eq:e7-schur-obstruction}
\end{equation}
If \(C\) remains a projector, \(c=0\) or \(1\), so its rank is zero or five,
never two.  Thus the exact branch projector must lose its gap, continuity, or
equivariance, or acquire new UV data.  This theorem does not prohibit every
non-idempotent equivariant coupling; it prohibits the exact separated
rank-two selector needed by the proposed light-only family.  A uniform exact
flavour action on the complete \(\bm5\) instead has
\begin{equation}
 \kappa_{\rm full}=2(+1)+2(-1)+0=0,\qquad
 \kappa_{\rm light}=+2,\qquad \kappa_{\rm heavy}=-2.
 \label{eq:e7-charge-trace}
\end{equation}
Therefore neither construction is yet an all-scale derivation of
\(k=+2\).

\subsection{Finite-site topology and the re-relaxation result}

The fixed-boundary finite-site configuration space of the declared naive
\(n+s\) action is
\begin{equation}
 {\cal C}_N=(S^3)^{(N-2)^3}\times\RR^{(N-2)^3}.
 \label{eq:e7-finite-config-space}
\end{equation}
Both factors are path-connected and a finite product of path-connected
spaces is path-connected.  Hence no continuous integer-valued function on
all of \({\cal C}_N\) can distinguish an exact \(B=1\) component from the
vacuum.  A finite-\(a\) topological sector requires an admissibility
restriction, topology-preserving interpolation, or infinite dislocation
barrier; it is not supplied by the naive sampled action.

For diagnostics, each cube is split into the six Freudenthal tetrahedra.  For
a generic target \(q\in S^3\), let
\begin{equation}
 V_T=[n_{x_0}\ n_{x_1}\ n_{x_2}\ n_{x_3}],\qquad
 c_T=V_T^{-1}q\quad\text{for }\det V_T\ne0.
\end{equation}
For a singular tetrahedron, \(\operatorname{span}V_T\) is a proper subspace
of \(\RR^4\).  The generic regular target is chosen outside the finite union
of these spans, so a singular cell has no preimage and contributes zero; an
exceptional target is equivalently handled by an arbitrarily small regular
perturbation.  The numerical implementation treats
\(|\det V_T|<10^{-13}\) as exceptional and separately verifies three-target
agreement and a nonzero-determinant margin for every counted hit.
The normalized affine tetrahedron hits \(q\) exactly when all entries of
\(c_T\) are positive, and its signed contribution is
\begin{equation}
 d_T(q)=\operatorname{sgn}(p)\operatorname{sgn}\det V_T,\qquad
 B_{\rm geom}=-\sum_Td_T(q).
 \label{eq:e7-geometric-degree}
\end{equation}
Three regular targets must agree.  The admissibility margin is
\begin{equation}
 \delta_a=\min_T\min_{\lambda_i\ge0,\,\sum_i\lambda_i=1}
 \left\|\sum_{i=0}^3\lambda_i n_{x_i}\right\|.
 \label{eq:e7-admissibility}
\end{equation}
Only on \(\delta_a>0\) is the normalized interpolation nonsingular and the
degree locally constant.

Five independent full-grid relaxations give:
\begin{center}
\begin{tabular}{rrrrrr}
\toprule
\(N\)&\(L\)&\(a\)&\(E_{\rm initial}\)&\(E_\star\)&\(B_{\rm geom}\)\\
\midrule
7&6&1.00&43.95516&\(9.07\times10^{-16}\)&0\\
9&6&0.75&54.15076&\(5.33\times10^{-17}\)&0\\
11&6&0.60&61.09804&\(8.07\times10^{-16}\)&0\\
9&8&1.00&43.16126&\(1.77\times10^{-14}\)&0\\
11&10&1.00&42.97910&\(2.04\times10^{-15}\)&0\\
\bottomrule
\end{tabular}
\end{center}
The maximum final free projected-gradient density is
\(1.84\times10^{-8}\): the optimizer has converged, but to the vacuum.
Pinning only \(n(0)=-1\) also gives \(B_{\rm geom}=0\) through a degenerate
tetrahedron.  Admissibility-guarded searches at
\(\delta_{\rm floor}=0.10,0.05,0.02\) retain \(B_{\rm geom}=1\), but stop on
the respective floor with gradient densities \(0.473,0.438,0.418\).  They do
not find an interior stationary point.

An explicitly artificial control fixes the central \(3\times3\times3\)
hedgehog \(n\)-values while re-relaxing all remaining variables independently
on each grid.  Its free constrained residual is below
\(6.2\times10^{-8}\), its three target degrees equal one, and its exact
restricted \(n+s\) Hessian has positive lowest eigenvalues
\(2.2643,1.8043,1.5002,1.1611,0.7563\) on the five grids.  The corresponding
full unanchored residuals are \(0.24\)--\(0.66\); therefore these are
restricted controls, not soliton stability eigenvalues.

On the same sites the audit also assembles a declared quadratic diquark
block, free Dirichlet gauge and ghost Laplacians, and
\begin{equation}
 K_F=D_W^\dagger D_W\succeq0.
 \label{eq:e7-fermion-normal}
\end{equation}
Their positive spectra are useful implementation checks.  However \(K_F\)
is not the second derivative of \(-\log\det D_W\), the interacting
gauge--meson cross blocks are absent, and a ghost operator is not an ordinary
positive bosonic Hessian block.  Consequently the direct sum is not a
physical BRST super-Hessian.  The exact finite-site conclusion is stronger
and more negative: the naive action has no exact topological component, and
its tested unconstrained minimizers unwind.

\subsection{The actual \texorpdfstring{\(SO(3)\)}{SO(3)} cover and the
rank-three alternative}

For the true localized orbit \(SO(3)\simeq\mathbb{RP}^3\), the cellular
boundary maps are \(0,2,0\).  Dualizing gives
\begin{equation}
 H^2(SO(3);\ZZ)=\ZZ_2=\{0,t\},\qquad t=\beta(a),\qquad 2t=0,
 \label{eq:e7-so3-cohomology}
\end{equation}
where \(a\in H^1(SO(3);\ZZ_2)\) is the cover class.  The unique nontrivial
complex line is the FR line
\begin{equation}
 {\cal L}_{\rm FR}=SU(2)\times_{\ZZ_2,\mathrm{sign}}\CC,\qquad
 (A,z)\sim(-A,-z).
 \label{eq:e7-fr-line}
\end{equation}
A generator of \(\pi_1(SO(3))\) lifts from \(A=\bm1\) to \(A=-\bm1\), so
\begin{equation}
 \operatorname{Hol}_\gamma({\cal L}_{\rm FR})=-1,\qquad
 \operatorname{Hol}_\gamma({\cal L}_{\rm FR}^{\otimes2})=+1.
 \label{eq:e7-fr-holonomy}
\end{equation}
For a real skew-adjoint family,
\begin{equation}
 \operatorname{Hol}_\gamma(\operatorname{Pf})
 =(-1)^{\operatorname{ind}_2({\cal D}_{Y_\gamma})},\qquad
 \operatorname{Det}\simeq\operatorname{Pf}^{\otimes2}.
 \label{eq:e7-pfaffian-line}
\end{equation}
Topology therefore classifies the two candidates but does not compute the
actual same-Yukawa mapping-torus mod-two index.  The standard conditional
Skyrmion rule gives
\begin{equation}
 (-1)^{N_cB}=(-1)^{2\cdot1}=+1,
 \label{eq:e7-fr-two-colour}
\end{equation}
and thus favours the trivial/bosonic sector, but it is not substituted for
the missing operator computation.  Inversion preserves the torsion class and
admits the topological Real lift
\begin{equation}
 J[A,z]=[A^{-1},\overline z],\qquad J^2=1,
 \label{eq:e7-fr-real-lift}
\end{equation}
while the microscopic CPT action on spin, gauge, regulator, domain, and
Pfaffian orientation remains open.

The torsion carrier cannot reproduce the old \(\cO(2)\) mechanism:
\(t\otimes\RR=0\) and a flat representative has zero curvature, whereas
\(c_1(\cO(2))\) is a nonzero free class.  Moreover the lowest odd FR rotor
block has \(j=1/2\) and full Peter--Weyl multiplicity four, not the three
holomorphic sections of \(\cO(2)\).

The independent rank-three branch does remove the AP-E6 two-band obstruction
at the bundle level.  On
\(X=\CP^1_z\times\CP^1_w\), set
\begin{equation}
 s=(z_0w_0^2,\ z_1w_1^2,\ z_0w_1^2+z_1w_0^2).
 \label{eq:e7-rank-three-sections}
\end{equation}
These three sections of \(\cO(1,2)\) have no common zero.  With
\begin{equation}
 e_+=\frac{\overline s}{\|s\|},\qquad
 P_+=e_+e_+^\dagger,\qquad
 {\cal M}=m(2P_+-\bm1_3),\quad m>0,
 \label{eq:e7-rank-three-mass}
\end{equation}
one has the exact global identities
\begin{align}
 \operatorname{spec}{\cal M}&=\{+m,-m,-m\},&{\cal M}^2&=m^2\bm1_3,
 \nonumber\\
 c_1(E_+)&=x+2y,&c_1(E_-)&=-(x+2y),&c_2(E_-)&=4xy.
 \label{eq:e7-rank-three-chern}
\end{align}
Thus \(E_+\oplus E_-=X\times\CC^3\), and rank three is minimal because the
\(\CP^2\) hyperplane class can have nonzero square while the \(\CP^1\)
two-band class cannot.  A scan of 11,664 points gives
\(\min\|s\|^2=0.25\), projector residual \(5.64\times10^{-16}\), spectral
residual \(2.89\times10^{-15}\), exact gap \(m=1.7\), and discrete Berry
Chern numbers \((1.000000000000003,1.999999999999998)\).

This is a genuine topological mass-family theorem, but its second
\(\CP^1_w\) factor is still a spectator rather than a localized normalizable
modulus.  No charged two-colour representation has yet produced its three
bands or the bidegree-\((1,2)\) Yukawa projector, and the finite matrix gap is
not a uniform Fredholm gap of the same soliton operator.

\subsection{AP-E7 decision}

Let \(C_{\rm APS}\) require the common regulator, the original target-family
lift, both actual target eta signs, the mixed heavy determinant, and all-scale
\(k=+2\).  Let \(C_{B1}\) require an unanchored admissible stationary sequence,
the interacting physical super-Hessian, and a dynamical continuum limit.
Let \(C_{\rm fam}\) require either the actual \(SO(3)\) Pfaffian route or a
physical rank-three same-soliton family, together with microscopic CPT and
gauge-basic WZW descent.  The current result is
\begin{equation}
 C_{\rm APS}=C_{B1}=C_{\rm fam}=\mathrm{false},\qquad
 \boxed{\text{degree-one Route-E portal not started}.}
 \label{eq:e7-final-decision}
\end{equation}
The mathematical advances are retained as reusable lemmas; they are not
promotion substitutes.  The canonical records are the three
\path{route_f/tex/ap_e7_*} notes, their machine-readable cards, and the
independent AP-E7 execution-frontier verifier.

\section{AP-E8 topology-preserving finite-grid \texorpdfstring{$B=1$}{B=1} construction}
\label{sec:ap-e8-topology}

Among the three proposed continuations, the topology-preserving action was
chosen because it directly changes the AP-E7 finite-site no-sector blocker.
The finite GW/domain-wall target kernel and the actual \(SO(3)\) Yukawa
mapping-torus index remain independent open lanes.

Let \(\mathcal E_F\) be the set of unique unordered vertex pairs that
co-occur in any tetrahedron of the fixed Freudenthal split.  For
\(-1/3<\epsilon<1\), \(\gamma>0\), and \(t>\epsilon\), define
\begin{equation}
 b_\epsilon(t)=-\log\frac{t-\epsilon}{1-\epsilon}
 +\frac{t-1}{1-\epsilon},\qquad
 E^{\rm TP}_{a}=E_a+\gamma a\sum_{\{x,y\}\in\mathcal E_F}
 b_\epsilon(n_x\cdot n_y),
 \label{eq:e8-action}
\end{equation}
and set the action to \(+\infty\) if any pair dot is at or below
\(\epsilon\).  The production values are \(\epsilon=0.01,\gamma=1\).
The subtracted linear term gives
\begin{equation}
 b_\epsilon(1)=b_\epsilon'(1)=0,\qquad
 b_\epsilon''(t)=\frac1{(t-\epsilon)^2},\qquad
 b_\epsilon(1-u)=\frac{u^2}{2(1-\epsilon)^2}+O(u^3).
 \label{eq:e8-barrier-calculus}
\end{equation}

For any Freudenthal tetrahedron whose six pair dots exceed \(\epsilon\),
every barycentric point obeys
\begin{equation}
 \left\|\sum_{i=0}^3\lambda_i n_i\right\|^2
 \ge\epsilon+(1-\epsilon)\sum_i\lambda_i^2
 \ge\frac{1+3\epsilon}{4}.
 \label{eq:e8-affine-bound}
\end{equation}
At \(\epsilon=0.01\), the normalized-affine norm is therefore at least
\(0.507444578255\).  The interpolation never meets zero, so its degree is
locally constant.  The divergent barrier prevents every continuous
finite-energy path from reaching the exceptional locus and changing \(B\).

There is also an exact finite-dimensional existence theorem.  In an energy
sublevel \(E_a^{\rm TP}\le C\), nonnegativity gives
\begin{equation}
 b_\epsilon(n_x\cdot n_y)\le\frac{C}{\gamma a}.
 \label{eq:e8-sublevel}
\end{equation}
Because \(b_\epsilon\) is strictly decreasing from infinity to zero, every
edge has a fixed-grid margin \(n_x\cdot n_y\ge\epsilon+\delta_{a,C}\).
The \(S^3\) variables are compact and the \(m_s>0\) potential is coercive in
all \(s_x\).  A minimizing sequence in any nonempty admissible component
therefore has an interior limit in the same component.  The limit is an
unanchored minimizer and all tangent-plus-scalar first variations vanish.
Sampling any smooth compact-supported continuum degree-one field makes the
\(B=1\) components nonempty on all sufficiently fine grids, hence supplies a
grid-indexed family of such minimizers.  No compactness uniform in \(a\) is
implied.

On a fixed uniformly smooth test field,
\begin{equation}
 n(x)\cdot n(x+av)=1-\frac{a^2}{2}|D_vn|^2+O(a^3),
\end{equation}
so the added action has the asymptotic form
\begin{equation}
 \gamma a\sum_e b_\epsilon(n_x\cdot n_y)
 =\frac{\gamma a^2}{8(1-\epsilon)^2}
 \int\sum_v\rho_v|D_vn|^4\,\dd^3x+O(a^3).
 \label{eq:e8-smooth-scaling}
\end{equation}
This estimate is not uniform on the unknown minimizers.  A compact-supported,
exact-boundary-compatible degree-one regression gives a global log--log slope
\(2.15910\), \(R^2=0.999141\), and successive local slopes
\(2.3356,2.1276,2.0695,2.0440\).

For tangent variations \(\xi_x\perp n_x\), the exact barrier second variation
on one edge, with \(q=n_x\cdot n_y\) and \(w=\gamma a\), is
\begin{align}
 \delta^2(wb_\epsilon(q))=w\{&b_\epsilon''(q)
 (\xi_x\cdot n_y+n_x\cdot\xi_y)^2\nonumber\\
 &+b_\epsilon'(q)[2\xi_x\cdot\xi_y
 -q(|\xi_x|^2+|\xi_y|^2)]\}.
 \label{eq:e8-edge-hessian}
\end{align}
The full Riemann Hessian must further contain the total sphere-curvature shift
\(Q^T[H_{\rm emb}-\operatorname{diag}(\lambda_xI_4,0)]Q\), with
\(\lambda_x=n_x\cdot G_x\).  AP-E8 validates
Eq.~\eqref{eq:e8-edge-hessian}; it does not yet assemble a physical aggregate
super-Hessian.

Seven independent final fields, with no core or centre pin, give
\begin{center}
\small
\begin{tabular}{rrrrrr}
\toprule
\(N\)&\(L\)&\(a\)&\(E_a^{\rm TP}\)&direct gradient density&\(\min n_x\cdot n_y\)\\
\midrule
15&6&0.428571&75.090097547&\(1.651\times10^{-7}\)&0.091845\\
17&6&0.375000&72.444298275&\(3.421\times10^{-7}\)&0.090442\\
19&6&0.333333&70.831336466&\(2.028\times10^{-7}\)&0.083178\\
21&6&0.300000&69.844122881&\(7.704\times10^{-7}\)&0.077786\\
16&6&0.400000&73.627393329&\(2.150\times10^{-7}\)&0.090722\\
21&8&0.400000&73.346860060&\(1.052\times10^{-7}\)&0.095418\\
26&10&0.400000&73.289478682&\(1.228\times10^{-7}\)&0.095975\\
\bottomrule
\end{tabular}
\end{center}
All three regular-value targets give \(B=(1,1,1)\) in every case.  The
fixed-\(a=0.4\) energy spread over \(L=6,8,10\) is \(0.459\%\), but the
weighted radius changes by about \(18\%\), so size/volume convergence is not
claimed.  Four nearby \((\gamma,\epsilon)\) controls and one deterministically
perturbed initial field also return admissible stationary \(B=1\) endpoints.
The total-gradient, invalid-continuation-gradient, and intrinsic-barrier-
Hessian residuals are \(9.38\times10^{-11}\), \(2.81\times10^{-11}\), and
\(8.16\times10^{-9}\).  The standalone card passes \(13/13\).

Thus the AP-E7 statement that no unanchored admissible finite-grid stationary
representative was available is superseded:
\begin{equation}
 \boxed{C_{B1}^{\rm finite\ grid}=\mathrm{true},\qquad
 C_{B1}^{\rm physical\ continuum}=\mathrm{false}.}
 \label{eq:e8-decision}
\end{equation}
Equicoercivity, Gamma convergence, a continuum stationary limit, barrier and
triangulation independence, the physical determinant Hessian, interacting
gauge/ghost blocks, four-dimensional QC2D dynamics, and the quantum continuum
limit remain open.  Hence the full Lane-B conjunction, portal authorization,
and physics promotion remain false.

\section{AP-E9: scaling regimes, Gamma-limit audit, and regulator no-go}

AP-E9 resolves ordinary compactness and the barrier scaling while retaining
the stronger topology and complete-action obligations.  Put
\begin{equation}
 d_a=1-\epsilon_a,\qquad w_a=\gamma_a a,\qquad
 R_a=\frac{\gamma_a a^2}{d_a^2}.
 \label{eq:e9-scales}
\end{equation}
For one admissible edge define
\begin{equation}
 x_e=\frac{n_x\cdot n_y-\epsilon_a}{d_a},\qquad
 z_e=1-x_e=\frac{|n_x-n_y|^2}{2d_a}.
\end{equation}
The subtracted logarithmic barrier is
$\phi(x_e)=-\log x_e-1+x_e$, with
\begin{equation}
 \phi(1-z)=\sum_{k\ge2}\frac{z^k}{k},\qquad
 \frac{z^2}{2}\le\phi(1-z)\le\frac{z^2}{2(1-z)}.
\end{equation}
Consequently
\begin{equation}
 \boxed{b_{\epsilon_a}(n_x\cdot n_y)
 \ge\frac{|n_x-n_y|^4}{8d_a^2}.}
 \label{eq:e9-quartic-bound}
\end{equation}

For every fixed $C>0$, if the total barrier is at most $C$, inversion by the
principal Lambert
branch gives the exact per-edge estimate
\begin{equation}
 x_e\ge -W_0(-e^{-1-C/w_a})
 \ge e^{-1-C/w_a}.
 \label{eq:e9-lambert}
\end{equation}
Thus uniform relative edge interiority on every positive bounded sublevel
holds exactly
when $\inf_a w_a>0$; a uniform absolute margin also requires
$\inf_a d_a>0$.  A single rotated vertex in a vacuum collar, with
$x_e=e^{-\alpha/w_a}$, gives a bounded-action countersequence when
$w_a\to0$.

Independently, the declared base coordinate Dirichlet term has a
mesh-uniform positive coefficient and gives an affine
$H^1$ bound on every fixed physical box.  Rellich compactness therefore
proves strong $L^2\times L^2$ equicoercivity for every positive $\gamma_a$
and $-1/3<\epsilon_a<1$.  This is not topological equicoercivity: strong
$L^2$ convergence does not preserve three-dimensional degree.

For a fixed smooth sphere-valued field and a periodic shape-regular mesh,
\begin{equation}
 \mathcal B_a[S_an]=\frac{R_a}{8}\int Q_{\mathcal T}(\nabla n)\,\dd x
 +O\!\left(R_a\left(a+\frac{a^2}{d_a}\right)\right).
 \label{eq:e9-smooth}
\end{equation}
At fixed $d$ and $\gamma_a=ca^{-p}$,
$w_a=ca^{1-p}$ and $R_a=ca^{2-p}/d^2$.  Uniform edge interiority needs
$p\ge1$, while smooth disappearance needs $p<2$; hence
\begin{equation}
 \boxed{1\le p<2,\qquad p=1\text{ is the minimal matched scaling}.}
 \label{eq:e9-window}
\end{equation}
If $d_a\sim d_0a^\alpha$, this becomes $p\ge1$ and
$p+2\alpha<2$; an absolute margin additionally rules out $\alpha>0$.

For the barrier alone, use the explicitly declared dual-cell
piecewise-constant embedding in $L^2(\Omega;\RR^4)$.  When $R_a\to0$, its
strong-$L^2$ Gamma limit is zero on $L^2(\Omega;S^3)$ and $+\infty$ outside;
it forgets trace and degree.  If $\inf R_a>0$, the quartic lower bound gives
$W^{1,4}$ and $C^{0,\beta}$ compactness, closing degree; if
$R_a\to\infty$, bounded sublevels converge to the vacuum.  At finite
nonzero $R$, the displayed $Q_{\mathcal T}$ is only the Cauchy--Born density
on fixed smooth samples.  A periodic multi-cell mesh generally requires a
homogenized cell formula, which is not evaluated here.  A resolved degree-one
bubble of radius
$r_a=\max\{R_a^{1/2},a^{1/2}\}$ has
$\mathcal B_a=O(R_a/r_a)\to0$ while converging in $L^2$ to the vacuum.
Therefore a continuum-vanishing barrier cannot by itself close degree.

At the critical scaling, the uncorrected smooth-sampling tensors are
algebraically different.  The uniform positive Freudenthal and checkerboard
five-tetrahedron quartics are
\begin{align}
 Q_{6,+}(A)&=\sum_i|Ae_i|^4+\sum_{i<j}|A(e_i+e_j)|^4
              +|A(e_1+e_2+e_3)|^4,\\
 Q_5(A)&=\sum_i|Ae_i|^4+\frac12\sum_{i<j}
 \left(|A(e_i+e_j)|^4+|A(e_i-e_j)|^4\right).
\end{align}
For $(Ae_1,Ae_2,Ae_3)=(v,0,0)$ their ratio is $4/3$, whereas for
$(v,v,0)$ it is $3$.  No scalar change of $\gamma$ equates these raw
Cauchy--Born tensors; equivalence of the homogenized densities remains an
open cell problem.

A complete same-action Gamma theorem is conditional: if the AP-E7 base
actions are equicoercive and Gamma-converge in strong $L^2\times L^2$, every
finite-energy limit has an energy-dense sequence of fixed smooth pairs, and
the nodal samples of each fixed pair recover the base energy, then a slow
diagonal argument together with
$a^2/d_a\to0$ and $R_a\to0$ imply that adding the barrier leaves that limit
unchanged.  The discrete forward-Skyrme-minor liminf, fixed-boundary
fixed-degree density, topological-current defect control, and translation
tightness for growing boxes remain unproved.

The production scan re-relaxes 23 fields over six fixed-box spacings, four
fixed-spacing volumes, $p=0,1,2$, both checkerboard phases, the original
six-tet mesh, and centered/translated starts at $a=0.4$ and $0.3$.  Every
case is stationary below $2\times10^{-6}$, admissible, and has three-target
$B=(1,1,1)$.  Nevertheless the predeclared continuum diagnostics are
\begin{center}
\begin{tabular}{lrrr}
\toprule
diagnostic&observed&threshold&pass\\
\midrule
relaxed $p=1$ barrier slope&0.2196&$1.0\pm0.4$&no\\
largest-three-volume radius spread&8.521\%&5\%&no\\
cross-mesh energy spread, $a=0.4$&5.442\%&3\%&no\\
cross-mesh energy spread, $a=0.3$&5.355\%&3\%&no\\
$p=0$ versus $p=1$ energy spread&6.446\%&3\%&no\\
\bottomrule
\end{tabular}
\end{center}
The energy-only volume spread is $0.1001\%$, demonstrating why energy alone
is an insufficient continuum gate.  The card passes $8/8$ implementation
and provenance checks but gives
\begin{equation}
 \begin{gathered}
 C_{B1}^{\rm physical\ continuum}=\mathrm{false},\\
 \texttt{regulator\_stable\_background}=\texttt{false},\\
 \texttt{same\_action\_Hessian\_allowed}=\texttt{false}.
 \end{gathered}
 \label{eq:e9-decision}
\end{equation}
The regulated determinant variation remains equally unauthorized.  The
finite GW/domain-wall kernel and actual $SO(3)$ Yukawa mapping-torus mod-two
index remain parallel lanes; the degree-one Route-E portal remains last.

\section{AP-E10: exact stencil no-go, finite-R cell formula, and translation quotient}

AP-E10 gives a decisive negative answer to the proposed compensated-
compactness route for the actual AP-E7 one-corner forward stencil.  On a
three-periodic torus let
\begin{equation}
 \theta_z=\frac{2\pi}{3}(z_1+2z_2),\qquad
 n_a(z)=\left(\sqrt{1-\eta^2a^2},
 \eta a\cos\theta_z,\eta a\sin\theta_z,0\right).
 \label{eq:e10-minor-sequence}
\end{equation}
Then $n_a\to e_0$ uniformly, but the target-component $(1,2)$ entry of the
spatial $(1,2)$ forward minor is identically
\begin{equation}
 \boxed{(D_1^an_a\wedge D_2^an_a)_{12}
 =\frac{3\sqrt3}{2}\eta^2.}
 \label{eq:e10-minor-defect}
\end{equation}
The continuum minor of the limit is zero.  Both the coordinate Dirichlet and
the AP-E7 forward Skyrme energy remain bounded.  Thus the minor weak-
continuity hypothesis needed in AP-E9 is false, not merely unproved.
This is consistent with div--curl theory requiring compatible differential
constraints \cite{BrianeCasadoDiaz2016} and with the FEEC requirement of a
discrete subcomplex and bounded cochain projection
\cite{ArnoldFalkWinther2010}.

A separate two-periodic construction tracks the current defect.  Put
\begin{equation}
 u_1=(-1)^{z_1+z_2},\quad u_2=(-1)^{z_2+z_3},\quad
 u_3=(-1)^{z_3+z_1},\qquad
 n_a=(\sqrt{1-3\eta^2a^2},\eta a u_1,\eta a u_2,\eta a u_3).
\end{equation}
Its image lies in an open hemisphere, so geometric degree is zero, whereas
the forward algebraic current is
\begin{equation}
 J_a=-16\eta^3\sqrt{1-3\eta^2a^2}\longrightarrow-16\eta^3.
 \label{eq:e10-current-defect}
\end{equation}
After multiplication by the exact-boundary cutoff
$\chi=\prod_i\sin^2(\pi x_i)$ and normalization, the weak limit is the
diffuse defect $-16\eta^3\chi^3\dd x$.  It is not a concentration measure.
Indeed, if $s_i$ are the three singular values of the forward derivative,
arithmetic--geometric mean gives
\begin{equation}
 \boxed{\int|J_a|^{4/3}\dd x\le
 \frac13\int|\wedge^2D^an_a|^2\dd x.}
 \label{eq:e10-l43}
\end{equation}
Thus singular forward-current concentration is excluded, but a diffuse
lattice alias survives and the forward current is not the normalized-affine
degree current.  Recent two-dimensional discrete Skyrmion compactness
\cite{BrianiCicaleseKreutz2026} supplies a useful architecture, not a theorem
for this noncompatible three-dimensional stencil.

The AP-E9 finite-$R$ cell problem is now solved.  For periodic bonds
$\mathcal B_Y$,
\begin{equation}
 Q_{\rm hom}(A)=\frac1{|Y|}\inf_{\varphi:Y\to\RR^m}
 \sum_{(p,r)\in\mathcal B_Y}|Ar+\varphi(p+r)-\varphi(p)|^4.
 \label{eq:e10-cell}
\end{equation}
Every fixed-direction bond graph in both the uniform six-tet and checkerboard
five-tet period-two cells is cycle-balanced.  Hence the mean periodic
corrector increment vanishes direction by direction; convex Jensen
inequality proves that $\varphi=0$ is the exact minimizer.  Therefore the
homogenized densities are precisely
\begin{align}
 Q_6^{\rm hom}(A)&=\sum_i|Ae_i|^4+\sum_{i<j}|A(e_i+e_j)|^4
 +|A(e_1+e_2+e_3)|^4,\\
 Q_5^{\rm hom}(A)&=\sum_i|Ae_i|^4+\frac12\sum_{i<j}
 \bigl(|A(e_i+e_j)|^4+|A(e_i-e_j)|^4\bigr).
\end{align}
The two checkerboard phases agree, but six/five ratios remain $4/3$ on
$(v,0,0)$ and $3$ on $(v,v,0)$.  Thus triangulation dependence survives
homogenization; a single rescaling of $R$ cannot remove it.  Random-start cell
solves reproduce the closed formulas to relative residual
$9.1\times10^{-15}$.

Finally, AP-E10 defines the translation quotient by the density
$\rho=1-n_0$, its barycentre, centered covariance, and barycentre-aligned
radial profile.  Dynamic shifts are accepted only if admissibility and all
three degree targets persist, after which the same action is re-relaxed.
Eleven production cases extend to $a=0.208333$ and $L=7.5$.  Every case is
stationary below $2\times10^{-6}$, has edge margin at least $0.1064$, and
gives $B=(1,1,1)$.  The predeclared continuum diagnostics are
\begin{center}
\begin{tabular}{lrrr}
\toprule
diagnostic&observed&threshold&pass\\
\midrule
joint-tail energy spread&1.656\%&3\%&yes\\
joint-tail centered-radius spread&9.612\%&5\%&no\\
joint successive profile $L^1$&7.072\%&5\%&no\\
translated-start profile $L^1$&3.500\%&3\%&no\\
five-/six-tet energy spread&5.143\%&3\%&no\\
five-/six-tet centered-radius spread&11.378\%&5\%&no\\
\bottomrule
\end{tabular}
\end{center}
The production card passes $10/10$, but the exact stencil counterexample and
failed mesh/shape gates force
\begin{equation}
 \boxed{\begin{gathered}
 \texttt{full\_Gamma\_liminf=false},\\
 \texttt{regulator\_stable\_background=false},\\
 \texttt{same\_action\_Hessian\_allowed=false}.
 \end{gathered}}
 \label{eq:e10-decision}
\end{equation}
The required repair is a complete tetrahedral cochain action: $dn$ as a
primal one-cochain, a graded shifted cup product satisfying a discrete
Leibniz rule, a positive calibrated Hodge star, and a topological three-
cochain whose evaluation agrees with normalized-affine degree.  The Hessian
and determinant variation remain embargoed until that replacement passes its
own cell and quotient continuum gates.

\section{AP-E11: compatible complete-minor cochain action}

AP-E11 replaces, rather than rescales, the AP-E10 stencil.  On an ordered
tetrahedral complex define the Alexander--Whitney front/back cup by
\begin{equation}
 (\alpha\smile_h\beta)[v_0\ldots v_{p+q}]
 =\alpha[v_0\ldots v_p]\,\beta[v_p\ldots v_{p+q}].
 \label{eq:e11-cup}
\end{equation}
The exact differential graded identities are
\begin{equation}
 d^2=0,\qquad d(\alpha\smile_h\beta)=d\alpha\smile_h\beta
 +(-1)^p\alpha\smile_h d\beta,
 \qquad (\alpha\smile_h\beta)\smile_h\gamma
 =\alpha\smile_h(\beta\smile_h\gamma).
 \label{eq:e11-dga}
\end{equation}
For $a^A=dn^A$, the normalized antisymmetric cochains
\begin{align}
 X_h^{AB}&=\frac12(a^A\smile_h a^B-a^B\smile_h a^A),\label{eq:e11-x}\\
 Z_h^{ABC}&=\frac1{3!}\sum_{\pi\in S_3}\operatorname{sgn}(\pi)
 a^{\pi(A)}\smile_h a^{\pi(B)}\smile_h a^{\pi(C)}
 \label{eq:e11-z}
\end{align}
are exactly the de Rham cochains of the P1 affine two- and three-minors.
The factor $1/2$ in \eqref{eq:e11-x} is the reference-triangle area.

For a $k$-simplex $\sigma=[i_0\ldots i_k]$, use the Whitney form
\begin{equation}
 \mathcal W_h(\sigma)=k!\sum_{r=0}^k(-1)^r\lambda_{i_r}
 d\lambda_{i_0}\wedge\cdots\widehat{d\lambda_{i_r}}\cdots
 \wedge d\lambda_{i_k}
 \label{eq:e11-whitney}
\end{equation}
and the consistent Hodge Gram matrix
$(M_k^T)_{\sigma\tau}=\int_T\mathcal W_h(\sigma)\wedge
\star\mathcal W_h(\tau)$.  Every $M_k^T$, $k=1,2,3$, is positive definite;
rigid rotations preserve it; and $\mathcal W_hd=d\mathcal W_h$ reconstructs
affine exterior forms exactly
\cite{Whitney1957,Dodziuk1976,ArnoldFalkWinther2010}.  The audit gives
$\lambda_{\min}=2.485\times10^{-2}$, rotation residual
$3.55\times10^{-15}$, and affine reconstruction residual
$2.40\times10^{-16}$.

The same third-cup family supplies topology.  On $T=[0123]$,
\begin{equation}
 C_h[n](T)=\epsilon_{ABCD}
 (n^A\smile_h a^B\smile_h a^C\smile_h a^D)(T)
 =\det[n_0\ n_1\ n_2\ n_3].
 \label{eq:e11-cup-det}
\end{equation}
If every tetrahedral vertex pair obeys $n_i\cdot n_j>\epsilon>-1/3$, then
for $v_T=\sum_i\lambda_i n_i$,
\begin{equation}
 |v_T|^2\ge\epsilon+(1-\epsilon)\sum_i\lambda_i^2
 \ge\frac{1+3\epsilon}{4}>0.
 \label{eq:e11-gap}
\end{equation}
Consequently $u_T=v_T/|v_T|$ is defined and its exact topological cochain is
\begin{equation}
 \Omega_h(T)=\frac{s_T}{2\pi^2}C_h[n](T)
 \int_{\Delta^3}|v_T(\lambda)|^{-4}\,d\lambda,\qquad
 \sum_T\Omega_h(T)=\deg u_h.
 \label{eq:e11-degree}
\end{equation}
This is the pullback formula for the $S^3$ volume form; the numerator is
constant on each affine tetrahedron and equals \eqref{eq:e11-cup-det}.

The compatible action is
\begin{equation}
 \begin{split}
 E_h&=\sum_T\int_T\left\{
 \frac12|Dv_h|^2+\frac R2|\wedge^2Dv_h|^2
 +\frac K2|\wedge^3Dv_h|^2+m^2(1-v_h^0)\right\}d^3x,\\
 &(R,K,m^2)=(1,0.35,0.12).
 \end{split}
 \label{eq:e11-action}
\end{equation}
The positive $K$ term is the Hodge norm of the complete family
\eqref{eq:e11-z}; its contraction with $n$ and radial dressing give
\eqref{eq:e11-degree}.  It closes the finite-grid topology-collapse mode
without reintroducing a mesh-direction quartic.

Every periodic P1 corrector preserves the means of all first, second, and
third minors.  Convex Jensen inequality on the complete minor vector proves
the exact cell formula
\begin{equation}
 \boxed{Q_{R,K}^{\rm cell}(A)
 =\frac12|A|^2+\frac R2|M_2(A)|^2+\frac K2|M_3(A)|^2,\qquad
 \varphi_{\min}=0.}
 \label{eq:e11-cell}
\end{equation}
It is identical on the uniform six-tet and both checkerboard five-tet meshes;
the optimized residual is $2.56\times10^{-14}$ and the minor-mean residual is
$1.78\times10^{-15}$.  This is the model-specific polyconvex mechanism
\cite{Ball1977}, not a scalar rescaling of AP-E10.

The Dirichlet term gives $H^1$ equicoercivity and forces the P1 limit back to
$S^3$.  The $L^2$ Hodge bounds on the compatible second and third minors,
together with Piola/null-Lagrangian identities, identify their weak limits.
Equation \eqref{eq:e11-gap} then identifies the normalized current
$\det[v_h,D_1v_h,D_2v_h,D_3v_h]/|v_h|^4$ and closes degree.  Convex lower
semicontinuity gives the liminf; smooth nodal interpolation gives recovery.
Thus the discrete theory Gamma-converges unconditionally to the fixed-degree
lower-semicontinuous relaxation.  Equality with the naive unrelaxed graph-
norm functional still requires a topology-sensitive smooth-density theorem
\cite{Bethuel1991}.

The production card passes $11/11$.  All thirteen unanchored cases have
$B=(1,1,1)$, projected gradient density at most $1.54\times10^{-7}$, and pair
dot at least $0.2724$.  The tail reaches $(N,L,a)=(37,8,0.222222)$; translation
and mesh comparisons use $(29,7,0.25)$.  With the same-action positive energy
density as quotient measure, the diagnostics are
\begin{center}
\begin{tabular}{lrrr}
\toprule
diagnostic&observed&threshold&pass\\
\midrule
joint-tail energy spread&0.437\%&3\%&yes\\
joint-tail centered-radius spread&3.489\%&5\%&yes\\
joint-tail radial CDF distance&4.067\%&5\%&yes\\
translated-start energy spread&$8.33\times10^{-12}$&1\%&yes\\
translated-start centered-radius spread&$7.16\times10^{-7}$&2\%&yes\\
five-/six-tet energy spread&0.463\%&3\%&yes\\
five-/six-tet centered-radius spread&0.752\%&5\%&yes\\
five-/six-tet radial CDF distance&3.556\%&5\%&yes\\
\bottomrule
\end{tabular}
\end{center}
Therefore
\begin{equation}
 \boxed{\begin{gathered}
 C_{\rm cell}=C_{\rm quotient}
 =C_{\rm relaxed\ regulator}=\mathrm{true},\\
 C_{\rm graph\ density}=C_{\rm classical\ regulator}
 =C_{\rm Hessian}=C_{\det}=\mathrm{false}.
 \end{gathered}}
 \label{eq:e11-decision}
\end{equation}
The AP-E10 stencil and mesh-universality blockers are closed.  The next
single mathematical gate is either fixed-degree smooth density in the
complete-minor graph norm or independent regularity and isolation of the
selected relaxed minimizer.  The classical same-action Hessian and regulated
determinant remain embargoed until one of those routes succeeds.


\section{AP-E12: endpoint graph density and relaxed-minimizer audit}

AP-E12 makes the AP-E11 regularity gate precise.  For fixed vacuum trace,
define
\begin{equation}
 \mathcal G_B^2=\{u\in W^{1,2}(\Omega;S^3):
 M_2(Du),M_3(Du)\in L^2,\ \deg u=B\}
 \label{eq:e12-graph}
\end{equation}
with strong $L^2$ convergence of $u$, $Du$, $M_2(Du)$, and $M_3(Du)$.
Graph convergence already closes degree: if $u_j,u$ are sphere-valued, then
\begin{equation}
 2\pi^2|B(u_j)-B(u)|\leq
 \|u_j-u\|_2\|M_3(Du_j)\|_2
 +|\Omega|^{1/2}\|M_3(Du_j)-M_3(Du)\|_2.
 \label{eq:e12-degree-continuity}
\end{equation}
Therefore a target-valued graph-dense smooth sequence is automatically in the
same integer degree for all sufficiently large indices.  The missing theorem
is target-valued endpoint graph density, not a separate degree repair.

For the AP-E11 continuum density
\begin{equation}
 W(F)=\frac12\{|F|^2+R|M_2(F)|^2+K|M_3(F)|^2\},
 \label{eq:e12-density}
\end{equation}
every minor of $F+t(a\otimes\xi)$ is affine in $t$.  Hence
\begin{equation}
 D^2W(F)[a\otimes\xi,a\otimes\xi]
 =|a|^2|\xi|^2+R|DM_2(F)[a\otimes\xi]|^2
 +K|DM_3(F)[a\otimes\xi]|^2,
 \label{eq:e12-lh}
\end{equation}
so the uniform Legendre--Hadamard constant is one.  More strongly, periodic
null-Lagrangian means give the exact identity
\begin{align}
 \int\{W(F+D\varphi)-W(F)\}
 =\frac12\int\{&|D\varphi|^2
 +R|M_2(F+D\varphi)-M_2(F)|^2\notag\\
 &+K|M_3(F+D\varphi)-M_3(F)|^2\}.
 \label{eq:e12-strong-qc}
\end{align}
This proves model-specific strong quasiconvexity but does not construct the
trace-preserving $S^3$ recovery sequence.

The classical Cartesian approximation result is exponent-losing:
$\operatorname{Cart}^{\bm p}$ is strongly approximable in every
$\mathcal A^{\bm q}$ with $\bm q<\bm p$, not at the endpoint
\cite{Maly1991}.  Area approximation controls all minors in $L^1$
\cite{Mucci2001}.  A scale-sharp audit also shows why the direct Mal\'y
two-hole dipole does not provide the missing counterexample here.  For
$u(r,\theta,z)=r^\alpha\gamma(\theta)$,
\begin{equation}
 \int|M_2(Du)|^2\sim\int r^{4\alpha-3}\,dr<\infty
 \quad\Longleftrightarrow\quad\alpha>\frac12,
 \label{eq:e12-dipole-energy}
\end{equation}
whereas Cauchy--Schwarz forces every core filling to pay at least
$c r^{4\alpha-2}$; a non-vanishing cost requires $\alpha\leq1/2$.  For
$r^{1/2}\log(e/r)^{-\beta}$, minor integrability requires $\beta>1/4$
while the filling bound tends to zero.  This excludes that mechanism, not
all possible endpoint concentrations.

The fallback cannot presently be closed by importing generic regularity.
The formal stress and classical equation are
\begin{align}
 P(F)&=F+R\{(\operatorname{tr}G)F-FG\}+KF\operatorname{cof}G,
 \qquad G=F^TF,\label{eq:e12-stress}\\
 (I-u\otimes u)&\{-\operatorname{div}P(Du)-m^2e_0\}=0.
 \label{eq:e12-el}
\end{align}
A relaxed minimizer is not authorized to satisfy \eqref{eq:e12-el} until a
local fixed-trace recovery theorem is supplied.  Moreover
$W(F)\lesssim1+|F|^6$ with $(n,p,q)=(3,2,6)$, while the 2024 smooth partial-
regularity theorem for relaxed strongly quasiconvex integrands requires
$q<\min\{np/(n-1),p+1\}=3$ \cite{GmeinederKristensen2024}.  It also does not
include the $S^3$ constraint and degree relaxation.

The production card passes $9/9$.  Three successively finer backgrounds have
maximum one-tetrahedron energy fractions
$3.966\times10^{-3}$, $1.849\times10^{-3}$, and
$1.173\times10^{-3}$.  Three amplitude-$0.24$ tangent perturbations re-relax
to energy spread $3.38\times10^{-11}$, maximum radial-CDF distance
$5.70\times10^{-7}$, and field RMS $1.31\times10^{-5}$ after quotienting
integer translations and proper target $SO(3)$.  This is finite-dimensional
same-basin evidence only, not continuum isolation.  The authorization is
\begin{equation}
 \boxed{\begin{gathered}
 C_{\rm graph\ degree}=C_{\rm strong\ QC}=C_{\rm dipole\ no\mbox{-}go}
 =\mathrm{true},\\
 C_{\rm endpoint\ density}=\mathrm{false\ (not\ disproved)},\\
 C_{\rm weak\ EL}=C_{\rm partial\ regularity}=C_{\rm classicality}
 =C_{\rm isolation}=C_{\rm Hessian}=C_{\det}=\mathrm{false}.
 \end{gathered}}
 \label{eq:e12-decision}
\end{equation}
The next theorem must be an annular target-valued endpoint replacement with
equiintegrable complete minors, or a reverse-H\"older estimate giving local
$L^{2+\delta}$ control of the complete-minor vector followed by a controlled
$S^3$ projection.  No same-action Riemann Hessian is assembled before one of
those paths also establishes classicality and continuum isolation.


\section{Reproducibility ledger}

The independent arithmetic/symbolic card is
\begin{center}
\path{route_f/code/verify_another_physics_route_e_bridge.py}.
\end{center}
Run from the repository root:
\begin{verbatim}
python3 route_f/code/verify_another_physics_route_e_bridge.py
\end{verbatim}
It writes
\begin{center}
\path{route_f/output/another_physics_route_e_bridge.json},\qquad
\path{route_f/output/another_physics_route_e_bridge.md}.
\end{center}
The current result is \(27/27\) mechanical checks, including
\begin{equation}
 \max|\Ktr^2-I/3|=1.11\times10^{-16},\qquad
 |2\Delta-B|=0,
\end{equation}
zero generator-invariance and canonical factor-two residuals, exact agreement
between the loaded Route-E \(\zeta\) card and the declared benchmark, complex
weight-two covariance error
\(2.80\times10^{-17}\), the two independent Q-ball energy estimates in
Sec.~6, and oscillatory-level-set residuals below \(2.7\times10^{-14}\).
The three earlier successor cards pass \(18/18\), \(38/38\), and \(27/27\),
respectively.  The new completion-frontier cards pass \(24/24\), \(58/58\),
and \(27/27\), while their integrated authorization gate passes \(20/20\).
The AP-E6 follow-up cards pass \(39/39\), \(20/20\), and \(25/25\);
their fresh-manifest, same-checksum frontier conjunction passes \(15/15\)
while keeping every complete-lane boolean false.
The AP-E7 regulator, discrete-topology, and \(SO(3)\)/rank-three cards pass
\(42/42\), \(16/16\), and \(18/18\), respectively.  Their independent
same-checksum execution-frontier conjunction passes \(16/16\), with the two
family alternatives represented by a disjunction but every condition within
each alternative retained conjunctively.  All complete-lane, portal, and
promotion booleans remain false.  AP-E8's topology-preserving finite-grid card
passes \(13/13\), promotes only the finite-grid sector/stationarity
subcondition, and explicitly keeps the physical-continuum Lane-B conjunction,
every complete lane, portal authorization, and promotion false.
AP-E9's scaling/triangulation card passes $8/8$, proves fixed-box strong-
$L^2$ equicoercivity, the barrier zero-limit, and the positive-$R$ compactness
regimes, while keeping the finite-$R$ homogenized density, numerical
continuum, same-action Hessian, and regulated-determinant gates false.
AP-E10's stencil/cell/quotient card passes $10/10$, explicitly disproves
one-corner forward-minor compensated compactness, excludes only singular
forward-current concentration, computes the finite-$R$ zero-corrector cell
densities, and proves their six-/five-tet inequivalence.  Its extended
stationary scan keeps the physical continuum, Hessian, and determinant gates
false.
AP-E11's compatible-cochain card passes $11/11$, constructs exact shifted-cup
Leibniz algebra and positive Whitney Hodge stars, identifies normalized-affine
degree with the shared third-cup contraction, proves the complete-minor five-/
six-tet cell formula, and passes every translation/triangulation quotient
threshold.  It promotes the relaxed regulator background only; the classical
graph-norm density, Hessian, and determinant gates remain false.
AP-E12's endpoint-density/regularity card passes $9/9$, proves graph-norm
continuity of degree, the exact strong-quasiconvexity identity, and a
scale-sharp no-go for the direct Mal\'y dipole mechanism.  The endpoint
density theorem is neither proved nor disproved.  The available relaxed
partial-regularity exponent range does not contain the present $(2,6)$
growth, so classicality, continuum isolation, Hessian, and determinant gates
remain false.
AP-E13's annular/reverse-H\"older card passes $9/9$ and disproves the
unconditional annular endpoint replacement by an explicit geodesic-cap
Wess--Zumino obstruction.  It proves a corrected replacement under
linear-in-radius target oscillation, while the strong-quasiconvexity
reverse-H\"older attempt remains circular at its highest cutoff-stress term.
Endpoint density itself is neither proved nor disproved, and all downstream
classicality, isolation, Hessian, determinant, and portal gates remain false.
The JSON explicitly retains all
cross-theory identifications as non-derived bridge axioms and sets
\texttt{physics\_promotion\_allowed=false}.

\section{AP-E13: annular endpoint replacement and reverse-H\"older audit}

AP-E12 proposed a target-valued annular replacement which keeps the trace on
a good sphere while controlling all three $L^2$ complete-minor families.
AP-E13 proves that its unconditional form is false.  For
\begin{equation}
 g_\varepsilon(\omega)=(\cos\varepsilon,\sin\varepsilon\,\omega)
 \in\mathbb S^3,
 \qquad
 u_{\rho,\varepsilon}(r\omega)=g_\varepsilon(\omega)
 \quad\hbox{on }A_{\rho,2\rho},
 \label{eq:e13-cap-trace}
\end{equation}
one has
\begin{equation}
 E_1=8\pi\rho\sin^2\varepsilon,\qquad
 E_2=\frac{2\pi}{\rho}\sin^4\varepsilon,\qquad E_3=0.
 \label{eq:e13-shell}
\end{equation}
The relative Wess--Zumino volume of the trace is
\begin{equation}
 V(\varepsilon)=2\pi\left(\varepsilon-\frac{1}{2}\sin2\varepsilon\right)
 =\frac{4\pi}{3}\varepsilon^3+O(\varepsilon^5)
 \quad\bmod 2\pi^2\mathbb Z.
 \label{eq:e13-wz}
\end{equation}
Gluing any trace-preserving filling $U:B_\rho\to\mathbb S^3$ to the canonical cap
and using integer degree gives
\begin{equation}
 \int_{B_\rho}|M_3(DU)|^2\dd x
 \geq\frac{V(\varepsilon)^2}{|B_\rho|}.
 \label{eq:e13-fill-lower}
\end{equation}
At $\varepsilon=\sqrt\rho$, the shell graph energy tends to zero but the
right side of \eqref{eq:e13-fill-lower} tends to $4\pi/3$.  Every sphere of
the shell has the same trace, so a good-radius choice cannot remove the
obstruction.  After zero extension to a fixed ball, the squared third-minor
densities retain positive mass on shrinking sets and are not equiintegrable.
This disproves the universal annular lemma, but not endpoint graph density
for every fixed map.

There is a scale-compatible repair.  If a trace $g:S_\rho\to\mathbb S^3$ lies in
a normal ball about $q$ and
$\|\log_qg\|_\infty\leq\Lambda\rho$, the geodesic cone
$U(r\omega)=\exp_q((r/\rho)\log_qg(\rho\omega))$ obeys
\begin{align}
 \int_{B_\rho}|DU|^2&\leq C(\rho T_1+\Lambda^2\rho^3),\notag\\
 \int_{B_\rho}|M_2(DU)|^2&\leq
 C(\rho T_2+\Lambda^2\rho T_1),\notag\\
 \int_{B_\rho}|M_3(DU)|^2&\leq C\Lambda^2\rho T_2,
 \label{eq:e13-conditional}
\end{align}
where $T_k$ are the tangential trace-minor norms.  Coarea on a good sphere
turns $\rho T_k$ into the corresponding shell cost.  Thus a linear Morrey
oscillation bound, or a comparably strong Wess--Zumino-flux decay theorem, is
the missing hypothesis.

The parallel reverse-H\"older route does not close from exact strong
quasiconvexity alone.  The classical stress
\begin{equation}
 P(F)=F+a_2\{(\operatorname{tr}F^TF)F-F(F^TF)\}
       +a_3F\operatorname{cof}(F^TF)
 \label{eq:e13-stress}
\end{equation}
contains a highest cutoff product
\begin{equation}
 \frac{|u-q|}{\rho}|M_2(Du)||M_3(Du)|
 \leq\eta|M_3(Du)|^2
 +\frac{C_\eta}{\rho^2}|u-q|^2|M_2(Du)|^2.
 \label{eq:e13-cutoff}
\end{equation}
The last term needs the same Morrey bound or prior higher integrability of
$M_2$, so using it to prove $L^{2+\delta}$ is circular.  The AP-E13 card
passes $9/9$: at $\rho=10^{-8}$ the critical shell norm is
$6.283\times10^{-8}$ while the filling lower bound is $4.188790188$.
Accordingly the conditional replacement theorem is true; the universal
replacement is false by counterexample; the model-specific reverse-H\"older
theorem, endpoint density, classicality, continuum isolation, Hessian,
determinant, and portal gates remain false.


\appendix

\section{Units and normalization checks}

In \(\hbar=c=1\) units, \([\Phi]=[f]=[m]=[M]=1\),
\([U]=4\), \([\sigma]=3\), \([R]=-1\), while \(Q,\lambda,\eta\) are
dimensionless.  Therefore every term in \eqref{eq:qballpotential} has mass
dimension four and every term in \eqref{eq:stepenergy} has dimension one.
For \eqref{eq:zetaportal}, \(g\) and \(\zeta_{\rm eff}\) are dimensionless
and \(\Lambda\) has dimension one.

The Q-ball charge convention differs by factors of two across the literature
depending on whether \(\Phi=\phi e^{\ii\omega t}\) or
\(\Phi=\phi e^{\ii\omega t}/\sqrt2\).  Equations
\eqref{eq:qcharge}--\eqref{eq:stepenergy} use one convention consistently;
mixing conventions would change \(Q\), \(f_0\), and \(R\) factors and is
forbidden in the verifier.

\section{Why a phase portal needs a second reference}

Suppose \(\zeta_{\rm eff}=c\Phi^2\) and define the putative invariant
\(I=\zeta_{\rm eff}^*\Phi^2/(|\zeta_{\rm eff}||\Phi|^2)\).  Then
\begin{equation}
 I=\frac{c^*\Phi^{*2}\Phi^2}{|c|\,|\Phi|^4}=\frac{c^*}{|c|},
\end{equation}
which is constant.  The field cannot serve as its own phase reference.  A
nontrivial observable requires a second spurion or matrix structure, such as
\(M_V\) in \(M_R=M_V+\zeta_{\rm eff}\Ktr\).  If that reference does not
transform, it explicitly breaks the Q-ball \(U(1)\); if it does transform, one
must solve the coupled background and identify an invariant relative phase.
This algebra is the sharp reason that phase doubling alone is not a visibility
mechanism.

\section{Minimal falsification criteria}

The bridge program should be abandoned or revised if any of the following
occurs:
\begin{enumerate}
\item the reduced bubble moduli space is not \(\CP^1\), or its physical bundle
has Chern number other than two;
\item the zero-mode operator has extra unpaired charged modes or does not
realize \(H^0(T)\);
\item the exact fixed-charge soliton has a negative Hessian mode or
\(E/Q\) above an allowed charge carrier;
\item every phase in \(M_V+\zeta\Ktr\) is removable, leaving no invariant
quantity such as \eqref{eq:phaseinvariants};
\item the required explicit breaking makes the soliton lifetime shorter than
the intended physical process;
\item a valid selection rule cannot forbid \(XLH\) and dangerous operators;
\item the embedded model fails Route-E flavor, threshold, proton-decay, or
cosmology gates;
\item a gravity carrier violates the redshift relation \eqref{eq:clockweak},
lensing, or fifth-force bounds.
\end{enumerate}

\bibliographystyle{unsrt}
\bibliography{route_f/tex/another_physics_route_e_bridge}

\end{document}
